% Chinese research blueprint; reported CPU results are scoped separately from proposed evaluations.
% !TeX program = xelatex
% Compile twice with XeLaTeX, or use Tectonic. The fallback is TeX Live's Arphic TTF.
\documentclass[11pt]{article}
\usepackage[a4paper,margin=0.82in]{geometry}
\usepackage{fontspec}
\usepackage{xeCJK}
% Use the bundled default Latin fonts; cover CJK roman, sans and mono separately.
\IfFontExistsTF{Noto Sans CJK SC}
 {\setCJKmainfont[AutoFakeBold=1.4,AutoFakeSlant=0.15]{Noto Sans CJK SC}
  \setCJKsansfont[AutoFakeBold=1.4,AutoFakeSlant=0.15]{Noto Sans CJK SC}
  \setCJKmonofont[AutoFakeBold=1.4]{Noto Sans CJK SC}}
 {% TrueType embeds ToUnicode without relying on an external Adobe-GB1 CMap.
  \setCJKmainfont[AutoFakeBold=1.4,AutoFakeSlant=0.15]{gbsn00lp.ttf}
  \setCJKsansfont[AutoFakeBold=1.4,AutoFakeSlant=0.15]{gbsn00lp.ttf}
  \setCJKmonofont[AutoFakeBold=1.4]{gbsn00lp.ttf}}
\usepackage{amsmath,amssymb,amsthm,mathtools}
\usepackage{booktabs,tabularx,longtable,array,enumitem,xcolor,hyperref,fancyhdr}
\usepackage{tikz}
\usetikzlibrary{positioning,arrows.meta}
\definecolor{navy}{RGB}{27,56,81}
\definecolor{lightblue}{RGB}{229,239,248}
\hypersetup{unicode=true,colorlinks=true,linkcolor=navy,citecolor=navy,urlcolor=navy,
 pdftitle={Intervention-Calibrated Distribution Transport for Cross-Embodiment Sim-to-Real Learning}}
\setlength{\parindent}{0pt}
\setlength{\parskip}{0.42em}
\setlength{\emergencystretch}{3em}
\renewcommand{\baselinestretch}{1.10}
\renewcommand{\arraystretch}{1.18}
\newcolumntype{Y}{>{\raggedright\arraybackslash}X}
\newcolumntype{L}[1]{>{\raggedright\arraybackslash}p{#1}}
\setlist{nosep,leftmargin=1.7em,topsep=0.25em}
\renewcommand{\contentsname}{目录}
\renewcommand{\abstractname}{摘要（研究设想）}
\renewcommand{\refname}{已核对的主要文献与实现入口}
\renewcommand{\figurename}{图}
\renewcommand{\tablename}{表}
\pagestyle{fancy}\fancyhf{}
\fancyhead[L]{\small 任务相关行为分布迁移}
\fancyhead[R]{\small 研究蓝图 / 2026-10-14}
\fancyfoot[C]{\thepage}
\setlength{\headheight}{18pt}
\newtheorem{proposition}{命题}[section]
\newtheorem{theorem}[proposition]{定理}
\newtheorem{assumption}[proposition]{假设}
\newtheorem{remark}[proposition]{说明}
\renewcommand{\proofname}{推导}
\newcommand{\E}{\mathbb E}
\newcommand{\R}{\mathbb R}
\newcommand{\Prob}{\mathbb P}
\newcommand{\ind}{\mathbf 1}
\newcommand{\cD}{\mathcal D}
\newcommand{\cQ}{\mathcal Q}
\newcommand{\cZ}{\mathcal Z}
\newcommand{\cI}{\mathcal I}
\newcommand{\cF}{\mathcal F}
\newcommand{\cP}{\mathcal P}
\newcommand{\cT}{\mathcal T}
\newcommand{\TV}{\operatorname{TV}}
\newcommand{\tr}{\operatorname{tr}}
\newcommand{\Lip}{\operatorname{Lip}}
\newcommand{\norm}[1]{\left\lVert #1\right\rVert}
\title{\vspace{-1em}\textbf{Intervention-Calibrated Distribution Transport\\for Cross-Embodiment Sim-to-Real Learning}\\[0.6em]
\large 干预校准的任务相关行为分布迁移\\[0.4em]
\normalsize 中文论文蓝图：方法为主，理论辅助}
\author{研究设计 / 修订版}
\date{2026 年 10 月 14 日}
\begin{document}
\maketitle
\vspace{-1em}
\begin{abstract}
不同机器人可以共享任务语义，但同一个技能分布在不同执行机制下会产生不同的接触、运动和失败分布。本文研究如何将源任务条件行为分布迁移为目标机器人可执行的行为分布。我们以已有源行为先验为起点，通过有限真实干预校准任务相关的执行机制，再学习一个 source-conditioned transport，使目标闭环响应兼顾任务完成、源行为结构和物理可行性。方法将任务调制的响应几何、受约束的机制校准、realization-aware transport 和摊销执行组织为一个框架；action decoding 仅承担身体特定的动作落地。理论用于解释相关响应的可辨识条件与校准误差对 transport 目标的影响。实验拟以全 3D、6/7 自由度机械臂为起点，依次研究持件进入受限工位、夹具辅助定位和带恢复的 keyed connector 插入；在 Panda 源臂、xArm7 工程桥接目标臂和UR5e 跨自由度目标臂上开展预算匹配验证。第四项实验保持同一插入任务和技能语义，将固定臂技能扩展到轮式移动机械臂、带臂四足及人形，检验有限校准能否处理移动身体与接触的耦合响应。实机型号与接口验证待确定，本文不预先声称已有部署或实验优势。
\end{abstract}
\textbf{本版主线。} 研究对象是\emph{条件行为分布及其闭环实现}。模拟样本和真实校准记录只是估计与优化的载体；源 replay 的筛选、加权 IQL 与样本量分配不再是方法主体。本文已有受限结构族的 CPU 机制研究与工程验证；完整工业任务、未知接触与实机结果尚待正式评估。

\textbf{结构。} Related Work 三组，Research Questions 三个；Section 3 为 Preliminary，且恰有三个子节。Method 四节，占主文约三分之一；理论主文保留两个辅助结果，其余推导进入附录。标题中的 distribution transport 不承诺普通 OT 或 Schr\"odinger bridge 的端点边缘分布假设。
\tableofcontents
\newpage

\section{Introduction}
\label{sec:intro}
\textbf{第一段：共享任务与不同的行为实现。} 不同机械臂可以完成同一定位或插入任务，却具有不同运动学、控制延迟、柔顺性和接触响应。一个源策略描述的不只是单条轨迹，而是对初态、扰动和多种动作选择的条件行为分布。将它部署到另一身体或真实系统后，实际产生的运动、接触和失败概率可能改变。因此，跨 embodiment 与 sim-to-real 可统一看作受目标执行机制约束的任务相关分布迁移问题。

\textbf{第二段：已有工作的三类进展。} 共同状态/动作表示帮助不同身体共享技能\cite{lsa,lad,xskill}；系统辨识与主动探索利用少量真实实验改善部署\cite{asid,spi,optimalexploration}；分布对齐方法已使用联合表示对齐与 OT/UOT 组织 sim-and-real 学习\cite{otreal}。Karen Liu 等的 Human2Sim2Robot 进一步展示了物体层任务语义在 embodiment gap 中的作用\cite{human2sim}。这些是本文的方法基础。

\textbf{第三段：本文集中解决的缺口。} 表示或指令分布接近，不自动意味着真实交互响应接近；辨识参数准确，也不自动决定源多模态行为应该如何在目标上改变。本文聚焦有限真实干预下，如何构造\emph{由目标机制实现、保留有效任务结构}的行为分布。目标成功轨迹分布通常没有被完整观测，因此它是由任务目标、源行为先验和校准目标机制共同决定的优化对象。

\textbf{可用于论文的一段英文 limitation。}
\begin{quote}
\raggedright
Existing cross-embodiment representations, sim-to-real adaptation, and distribution-alignment methods provide important tools for sharing behavior across domains, but aligning nominal representations or commands does not by itself establish how the resulting behavior distribution will be realized under a different execution mechanism. We study how limited task-relevant interventions can calibrate this mechanism and guide the transport of a source behavior distribution into a feasible target distribution while retaining useful task structure.
\end{quote}

\textbf{第四段：具体问题。} 两种机械臂可能具有相近的自由空间轨迹，却在孔口侧向偏移后产生恢复与卡滞两种不同分布。源先验中的多个绕行路径或接触调整方向，也未必全部适用于目标工作空间。迁移需要改变目标行为的实现和概率结构，并允许必要的任务相关偏离；无条件追求两个分布越接近越好并不合理。

\textbf{第五段：方法概览。} 从已有源先验采样短时技能，通过少量真实探针校准目标响应，以任务阶段相关的几何定义哪些差异重要，再学习 source-conditioned transport。确定性映射已能保留随机源 prior 的多峰性；只有需要表达额外条件不确定性时才引入 transport noise，并单独验证其收益。训练在校准模型中进行，部署根据当前实际状态直接生成目标技能并闭环执行。action decoding 使用现有机器人接口。

\subsection{Three Research Questions}
\begin{enumerate}[label=\textbf{RQ\arabic*.}]
\item \textbf{Transfer effectiveness.} 在相同目标交互预算下，迁移后的行为分布能否改善跨机械臂与 sim-to-real 的闭环任务表现，并将同一操作技能扩展到移动身体？
\item \textbf{Calibration efficiency.} 有限任务相关干预能否校准影响迁移的执行响应，并在已声明机制族中泛化到未见初态、控制器和接触条件？
\item \textbf{Transport mechanism.} 干预响应、任务几何与可行性约束分别如何影响目标分布的任务对应、多模态覆盖和执行稳定性？
\end{enumerate}

\subsection{Three Candidate Contributions}
\begin{enumerate}[label=\textbf{C\arabic*.}]
\item \textbf{可执行分布迁移的建模。} 将任务条件源行为先验与目标闭环行为分布分开，通过目标执行机制定义 transport 的可实现集合；目标 marginal 是待构造的行为 law。
\item \textbf{干预校准的 realization-aware transport 方法。} 结合任务调制响应几何、低维机制校准、可执行约束和 source-conditioned mapping，学习目标行为分布；映射是否需要额外随机性由消融决定，action decoding 为执行接口，已有生成模型可直接作为 prior。
\item \textbf{预算匹配的分布与任务验证。} 在跨机械臂、接触操作与真实平台上，同时检验闭环效果、干预预算、多模态覆盖和响应分布变化，并用辅助分析解释适用条件。
\end{enumerate}
主要贡献由方法与实验建立。任务相关探索、OT/UOT、flow matching、Lagrange 更新与动作映射均有前例，不单独宣称为新概念。

\section{Related Work}
\label{sec:related}
\subsection{Cross-Embodiment Representations and Behavior Transfer}
Latent Space Alignment 使用编码/解码、cycle consistency 与 latent dynamics 对齐跨机器人技能\cite{lsa}；Latent Action Diffusion 以语义对齐 latent action 支持不同末端执行器的 diffusion policy\cite{lad}；XSkill 研究跨 embodiment 技能发现\cite{xskill}。Human2Sim2Robot 使用物体轨迹构造共同任务奖励\cite{human2sim}。本文继承共同任务语义，关注给定表示后，行为分布在不同机制下的闭环实现。

\subsection{Sim-to-Real Adaptation and Active System Identification}
Dynamics randomization、GAT 和 RMA 分别提供随机化、动作 grounding 与历史适应路线\cite{dr,gat,rma}。ASID 与 SPI-Active 主动采集信息并辨识目标机制\cite{asid,spi}；Guanya Shi 参与的 optimal exploration 已研究任务相关实验设计\cite{optimalexploration}。本文把有限实验得到的机制信息用于构造目标行为 law，而不将 task-aware exploration 本身作为主要创新。

\subsection{Distribution Alignment and Transport for Robot Learning}
OT-sim2real 通过联合表示分布对齐、UOT 与时间相关采样进行 sim-and-real co-training，其主要协议使用少量目标演示\cite{otreal}。DARA 与 H2O+ 研究动力学差异下的源/目标学习\cite{dara,h2oplus}。Flow matching 的路径成本和约束结合也已有研究\cite{densityflow}。本文重点在有限干预校准的接触响应、目标可执行性与源行为结构之间的联合建模；不能把联合分布对齐、partial transport 或受约束 flow 单独列为新贡献。

\section{Preliminary}
\label{sec:prelim}
\subsection{Embodiment-Conditioned Control and Action Decoding}
环境 $e$ 有自己的原始观测 $o_e$、动作 $a_e$ 和低层闭环控制。共同状态 $q=\phi_e(h_e)$ 包括物体/TCP 相对位姿、任务阶段、接触与力觉历史；必要时加入机器人能力和配置摘要。不同身体的关节维度不强行混合。

共同技能指令 $U\in\R^{H\times d}$ 表示短时任务空间命令、力设定值及允许的顺应参数。通过已有 action decoding 接口 $\mathcal A_e$ 与反馈控制器 $C_e$ 执行：
\begin{equation}
 a_{t+j}=C_e\big(h_{t+j};\mathcal A_e(U_j,q,c_e)\big),
 \qquad j=0,\ldots,H-1.
 \label{eq:decode}
\end{equation}
编码坐标系、单位、控制时长和允许动作边界需统一。实际接触力为响应，force setpoint 为可设置指令。这里不提出新 decoder 架构，也不把 decoder 单列为创新模块。理论假设共同状态在声明支持内足以描述宏步交互。

移动身体沿用任务坐标系中的同一技能 $U$；$\mathcal A_e,C_e$ 负责将末端参考落地到机械臂、底盘或腿部的全身控制。目标身体状态另包含基座位姿/速度、支撑接触、关节余量与控制器历史，不能只输入 TCP 位姿就假定跨身体状态充分。共享响应对应只比较任务进展、工件/TCP 与任务接触；身体专属的防倾覆、足端/轮地滑移及执行边界进入目标可行性与风险，不要求移动身体复制固定臂的基座运动。已有 UMI-on-Legs 采用 task-frame 末端轨迹连接操作策略与全身控制，并展示固定臂策略向四足平台的迁移；本文需在这一接口基础上验证校准与 transport 的增量价值，而非将接口本身作为新贡献\cite{umilegs}。

\subsection{Conditional Behavior Distributions and Transport}
源 prior 是已有源策略或生成模型：
\begin{equation}
 U_s\sim p_s(U\mid q,g).
\end{equation}
学习 transport kernel $T_\varphi(dU_t\mid U_s,q,g,c_t)$，诱导目标指令分布
\begin{equation}
 p_{t,\varphi}(dU_t\mid q,g)
 =\int T_\varphi(dU_t\mid U_s,q,g,c_t)\,p_s(dU_s\mid q,g).
 \label{eq:kerneltransport}
\end{equation}
确定性特例为 $p_{t,\varphi}=\Phi_\varphi\#p_s$。目标实际行为 law 为
\begin{equation}
 \mathbb P_{t,\varphi}^{g}
 =\operatorname{Rollout}_{K_t,\pi_{t,\varphi}}\#\mu_0.
 \label{eq:pathlaw}
\end{equation}
$p_{t,\varphi}$ 与 $\mathbb P_{t,\varphi}^{g}$ 是不同对象：前者为指令层，后者由目标反馈、动力学和扰动共同实现。源与目标响应使用共同初态/任务条件比较；不把不同初态的边缘分布当作机制差异。

样本重加权可以改变已有支持上的质量；transport 还可以改变目标指令的支持位置。一个可解析例子是 $y_s=u_s$、$y_t=\kappa u_t$：源 prior 位于 $[-1.1,-0.9]\cup[0.9,1.1]$，$\kappa=0.5$ 时，$u_t=2u_s$ 才能保持对应输出。仅在原 compact support 内改权重无法产生该指令分布。若目标边界禁止 $2u_s$，则必须接受不可消除的任务响应偏离。

\subsection{Problem Formulation}
输入包括源行为 prior、共同任务目标、可用目标机器人模型/控制接口，以及总预算 $B$ 的真实校准交互。输出为满足执行约束的目标条件行为分布。目标成功轨迹 law 不预先给定。

给定当前 $q,g$，以固定长度 $H$ 的条件轨迹 law 建立短时 formulation：
\begin{equation}
 \begin{aligned}
 \min_{Q,\Gamma}\quad&
 \E_{\Gamma}\big[C_g(\tau_t)+\lambda c_g(\tau_s,\tau_t)\big],\\
 \text{s.t.}\quad&\Gamma_s=\mathbb P_s^{g,H}(\cdot\mid q),\quad \Gamma_t=Q,
 \quad Q\in\mathfrak P_t^H(q),\\
 &\text{目标执行与风险约束},
 \end{aligned}
 \label{eq:ideal}
\end{equation}
$\mathfrak P_t^H(q)$ 为目标机制及允许短时闭环执行诱导的分布集合。实现中用 $T_\varphi$ 参数化其中的子集，并由源条件与随机种子构造 coupling。部署在每个目标实际状态重新条件化，诱导 Eq.~\eqref{eq:pathlaw} 的全任务 law；拼接的源参考一般不保持独立源 rollout 的全任务 marginal。本文不声称求得所有 coupling 的全局 OT 最优解。

主协议 A 仅使用声明校准实验，不使用目标任务演示或目标 RL fine-tuning。附加协议 B 给所有方法相同目标任务缓冲区；需要目标演示的 OT-sim2real、LAD 等在其适用监督协议下比较。所有校准、动作接口拟合和验证计入 $B$；最终评估不参与训练。

\section{Method: Intervention-Calibrated Distribution Transport}
\label{sec:method}
\textbf{核心流程。} 固定源 prior，建立任务响应表征；用有限可行干预校准目标机制；在校准模型中学习 transport；部署时按当前实际状态采样并执行目标技能。method 的新意需要来自这四步的配合及实证效果。

\begin{figure}[ht]
\centering
\begin{tikzpicture}[node distance=0.6cm and 0.6cm,>=Stealth,
 box/.style={draw=navy,rounded corners,fill=lightblue,align=center,font=\small,
 text width=3.65cm,minimum height=1.02cm}]
\node[box] (prior) {源条件行为分布\\已有策略 / diffusion / flow prior};
\node[box,right=of prior] (cal) {有限真实干预\\校准目标响应与不确定性};
\node[box,below=of prior] (geom) {任务响应几何\\进展、接触与模式变化};
\node[box,below=of cal] (trans) {学习 stochastic transport\\目标机制、任务与可行约束};
\node[box,below=of trans] (exec) {目标行为分布\\action decoding 与闭环执行};
\draw[->] (prior)--(trans);
\draw[->] (cal)--(trans);
\draw[->] (geom)--(trans);
\draw[->] (trans)--(exec);
\draw[->] (exec.west) -- ++(-0.3,0) |- (prior.south);
\end{tikzpicture}
\caption{方法关系。校准模型用于训练 transport；部署依据实际反馈重新条件化源 prior。回路中的重新采样不是在线轨迹 refinement。}
\end{figure}

\subsection{Task-Modulated Response Representation}
\label{sec:geometry}
\textbf{表征对象。} 每个技能执行产生响应描述 $Y=\Psi_g(\tau)$：任务进度、物体/TCP 几何、接触力统计与时间序列、跟踪误差、接触模式及失败事件。源记录的未来轨迹仅用于离线训练；在线 transport 输入只有当前可用历史。

\textbf{任务几何。} 使用有量纲归一化的距离
\begin{equation}
 d_g^2(y,y';q)=(y-y')^\top G_g(q)(y-y'),\qquad
 G_g(q)=\epsilon I+\sum_j\omega_j(q,g)b_jb_j^\top,\quad\epsilon>0.
 \label{eq:metric}
\end{equation}
$\omega_j$ 由阶段、接触状态、目标区域和约束活跃程度指定。受限接近突出全臂/持件间隙与姿态进展，夹具定位突出压靠、滑移与旋转；插入突出横向接触、推进和卡滞/恢复分支。模式概率差异通过 one-hot/事件成本独立进入，不能用平均力掩盖不同失败概率。$\epsilon I$ 避免忽略将影响未来决策的方向；若只保留部分特征，则理论与主张限制在该 quotient space。

\textbf{轨迹成本。} 对相近任务阶段的轨迹定义
\begin{equation}
 c_g(\tau_s,\tau_t)=
 \sum_{k=0}^{H}d_g^2(\Psi_{g,k}(\tau_s),\Psi_{g,k}(\tau_t);q_k^s)
 +\lambda_m c_{\rm mode}(\tau_s,\tau_t).
 \label{eq:semantic}
\end{equation}
首版固定技能时长；明显速度不同的机制可采用受限 phase alignment，并单独保留完成时间成本。不能用任意 time warping 消除真实延迟和力过冲。共同任务成本 $C_g$ 依据进展、接触保持、失败与完成定义；源语义成本与目标任务成本同时存在，允许目标改变源行为以完成任务。

\textbf{多模态结构。} 在 source prior 中保留多个有效模式 $m$，如左右调整或不同绕行路径。第一版用可解释几何规则标记模式，并在目标可执行的模式集合内测量覆盖。mode 标签帮助 transport 避免把多种行为平均成一个无效指令；它不保证所有源模式在目标上都应保留。

\textbf{实现与消融。} 表征使用 TCP/物体状态与 wrist F/T，先不引入大型视觉模型。比较 uniform metric、仅名义几何、平均响应和包含模式/扰动响应的完整表征。这部分吸收 2GO 的任务调制几何 insight。

\subsection{Budgeted Intervention and Target-Mechanism Calibration}
\label{sec:calibration}
\textbf{机制模型。} 在多机制模拟器中预训练短时响应模型
\begin{equation}
 (q',Y)\sim\widehat K_\psi(\cdot\mid q,U,\delta,c),
 \label{eq:model}
\end{equation}
$c$ 为低维执行 context，例如指令增益、延迟、接触顺应性或摩擦参数。真实交互校准 $c_t$ 及预测不确定性：
\begin{equation}
 \hat c_t=\arg\min_c
 \sum_{(q,U,\delta,y)\in\mathcal B_{\rm fit}}
 -\log\widehat K_\psi(y\mid q,U,\delta,c).
 \label{eq:calibrate}
\end{equation}
已有目标模拟器可先校准参数，再训练可微的短时 surrogate。所有目标真实交互（含进入接触阶段、失败和独立验证）计入 $B$，并同时报告控制步数与交互秒数；公开模拟预训练和 calibrated-twin 优化另列模拟计算预算，所有对照具有相同的可用模型与计算额度。少量真实样本用于校准结构化模型，不从头学习任意 world model。物理参数 context 必须限制在声明的物理支持内；离散 delay 用离散假设比较。自由 latent 的坐标不能继续解释为物理参数真值。

\textbf{干预语义。} $\delta$ 是可随机分配的命令、初始偏移或实验配置。定义规定执行接口下的响应
\begin{equation}
 K_e^\delta(dy\mid q,U)
 =\Prob_e(Y\in dy\mid q,\operatorname{do}(U),
 \operatorname{do}(\Delta=\delta);\mathcal A_e,C_e).
 \label{eq:intervention}
\end{equation}
使用匹配初态、已记录历史和随机指派，减少行为选择偏差。第一版探针为自由空间增益/延迟、接触建立、不同切向速度、横向偏移与受限恢复。力监测使用已有执行阈值，不从少量数据推导未知接触的全局安全性。

\textbf{预算选择。} 设探针为有限候选集 $\Delta_F$，每个候选满足 MDOC 式的已知可行性约束。第一版使用已有 task-dependent design，在模型中评估相关响应不确定性的下降：
\begin{equation}
 \min_{\mu}\ \Phi_g(\mu),\qquad
 \E_\mu c(\delta)\le b,\quad\operatorname{supp}\mu\subseteq\Delta_F,
\end{equation}
\begin{equation}
 \Phi_g(\mu)=
 \begin{cases}
 \tr(B_gF_\mu^\dagger B_g^\top),&\ker F_\mu\subseteq\ker B_g,\\
 +\infty,&\text{其他情况}.
 \end{cases}
 \label{eq:probe}
\end{equation}
这是方法支撑，非首次提出任务相关探索。增强版可用 ensemble 的虚拟测量更新，选择降低 transport 成本不确定性最多的探针；它应单独消融，不能与固定设计 Gaussian 公式混用。

\textbf{校准与验证分开。} 拟合集用于机制/接口参数，独立验证集检查各任务模式、响应分布与不确定性。动作接口调整后重新验证；所有拟合、验证和失败探针均计入 $B$。未覆盖的接触分支作为 unknown 报告。source prior 或 transport 改变访问分布时，重新检查模型支持，避免把校准工装覆盖误当作任意任务覆盖。

\subsection{Realization-Aware Distribution Transport}
\label{sec:transport}
\textbf{参数化。} 最小原型为 source-conditioned residual stochastic map：
\begin{equation}
 U_t=\Phi_\varphi(U_s,q,g,\hat c_t,m,\varepsilon)
 =\operatorname{Bound}\big(U_s+\Delta_\varphi(U_s,q,g,\hat c_t,m,\varepsilon)\big).
 \label{eq:residual}
\end{equation}
$\varepsilon$ 提供随机性，$m$ 表示源行为模式，$\operatorname{Bound}$ 实施已知命令边界。真实目标响应由 $K_t$ 实现，训练预测由 $\widehat K_\psi$ 实现。约束与损失在最终执行指令上评估，不能仅检查投影前输出。

\textbf{反馈与表达支持。} 源行为实现包含实际接触阶段与位姿误差；阶段由接触、撤离和对准事件推进，不能只用绝对时钟替代。受限 CPU prior 只预测 phase/mode 条件 reference，位姿误差反馈由所有方法共享的 decoder 实现；更丰富 prior 可直接以这些反馈条件化。源停止不意味着目标已到平衡点，因此不得用 $|U_s|$ 将目标所有对应纠偏方向硬置零。采用有界残差及实际 $SE(3)$ 反馈维持终止平衡；力控制期间投影法向位置参考，避免未使用参考积累后在释放时产生突变。反馈 decoder 对所有 baseline 相同，其修正不作为 transport 的独立贡献。

若 residual map 有效，再使用 conditional flow 增强表达：
\begin{equation}
 \frac{dX_\ell}{d\ell}=v_\varphi(X_\ell,\ell;U_s,q,g,\hat c_t,m),
 \quad X_0\sim\mathcal N(0,I),\quad U_t=\operatorname{Bound}(X_1).
 \label{eq:flow}
\end{equation}
这是一个 source-conditioned transport kernel；$\ell$ 是生成/transport 时间，不是物理时间。基础 residual map 与 flow 必须使用相同任务目标比较，不把新的 diffusion/flow 架构作为必需贡献。

\textbf{实际响应目标。} 对 source sample 与 transported sample 在匹配 $q,g,\delta$ 下生成源/目标预测轨迹。主损失为
\begin{equation}
 \begin{aligned}
 J(\varphi)=\E\big[&C_g(\hat\tau_t^\varphi)
 +\lambda_s c_g(\tau_s,\hat\tau_t^\varphi)
 +\lambda_u\|U_t-U_s\|^2_{H_g}\big]
 +\lambda_M\mathcal L_{\rm mode}(\varphi).
 \end{aligned}
 \label{eq:objective}
\end{equation}
第一项驱动目标任务完成；第二项保留实际任务结构；第三项限制无必要的指令改变。$H_g$ 是具有单位归一化的指令度量，不能直接把响应维度的 $G_g$ 乘到不同维度的指令上。可选地用模型 Jacobian 构造 $H_g=J_\psi^\top G_gJ_\psi+\epsilon I$，并单独检查其计算代价。

$\mathcal L_{\rm mode}$ 使用冻结的几何规则/分类器，只对已知几何与预算内校准模型支持的源模式约束输出模式信息；不读取最终诊断 oracle 的目标真值可行性标签。不能训练一个自由分类器通过记忆 source label 伪造模式保留。模型判为不支持或未知的模式不强制匹配，实际输出的改变、停止或恢复必须进入整体任务评估。

\textbf{从 coupling 到 distribution。} Eq.~\eqref{eq:objective} 对 $U_s\sim p_s$ 与 $U_t\sim T_\varphi(\cdot\mid U_s)$ 的联合 law 优化，其目标 marginal 是 Eq.~\eqref{eq:kerneltransport}。因此它在改变整个目标行为分布及其支持，而非给源 replay 赋权。需要比较的 realized response law 为
\begin{equation}
 \rho_{t,\varphi}^{\delta}(dy\mid q,g)
 =\int K_t^\delta(dy\mid q,U_t)\,
 T_\varphi(dU_t\mid U_s,q,g,\hat c_t)\,p_s(dU_s\mid q,g).
 \label{eq:responselaw}
\end{equation}
为避免 samplewise coupling 过强，可增加批内 Wasserstein/UOT 响应对齐，但它是已有工具的增强版。基准实现先使用可解释的 coupling 和目标任务成本，再比较 nominal feature UOT 与 realized-response UOT。UOT 中未匹配质量不是消失的实际任务；实际策略仍是归一化行为 law。

\textbf{硬约束与接触风险。} 已知运动学、速度/命令变化率和工作空间边界进入可行域。模型相关风险使用 ensemble：
\begin{equation}
 \mathcal R_j(\varphi)=
 \sup_{M\in\mathcal M_{\mathcal B}}
 \operatorname{CVaR}_{1-\delta_j}
 \left[\max_k h_j(\hat q_k,\hat U_k)\right]\le b_j.
 \label{eq:risk}
\end{equation}
必要时用独立误差估计 tightening。该风险是模型条件下的约束，不是实机安全证书；真实执行仍遵守已有验证范围与监测。

\textbf{受约束优化。} 把 MD-COAS 的 hard/soft 协调 insight 具体化：
\begin{equation}
 \mathcal L(\varphi,\lambda)
 =J(\varphi)+\sum_j\lambda_j(\mathcal R_j(\varphi)-b_j),
 \qquad\lambda_j\ge0,
\end{equation}
\begin{equation}
 \varphi\leftarrow\varphi-\eta_\varphi\nabla_\varphi\mathcal L,
 \qquad
 \lambda_j\leftarrow[\lambda_j+\eta_\lambda(\widehat{\mathcal R}_j-b_j)]_+.
 \label{eq:lagrange}
\end{equation}
hard constraints 通过参数化/可行支持实施，soft semantic costs 通过目标体现，风险约束由乘子协调。此结构用于离线 transport 学习，不重新实现 MGA 的在线 annealing/LRC。

\textbf{训练实现。} 首版采用短时可微 surrogate，固定 $p_s$ 和校准模型，对 transport 反向传播。模型先以公开训练机制预训练，再用目标 context 校准；隐藏目标真值不参与训练。若可实例化校准 physics twin，可在其中生成同预算数据或用已有 SAC 学习 residual actor，并保持源语义与风险目标。仅有 latent surrogate 时，换优化器仍受同一模型误差限制。两条路线分别报告；不声称黑箱接触天然可微。

\textbf{本轮可复核的有限实现。} 当前 P1.6/P1.7 使用结构化 calibrated twin
和冻结的 phase/mode reference prior，以有限候选枚举代替上述可微 learner。
在线反馈由所有方法共享的 decoder 承担；候选改变阶段阻抗、终段力和切向偏置，
P1.7 可逐模式选择相同族中的映射。先做有限 reset 的经验可行性过滤，再按
Section~\ref{sec:finite-feasibility} 的任务预算带选择响应对应；支持不足返回拒绝。
这条机制路线不等同于已完成 neural response、完整语义 source query 或大型
随机 transport。它与同类 task-only 具有相同表达类，额外价值须由冻结结果建立。

\subsection{Training Algorithm and Amortized Closed-Loop Deployment}
\label{sec:deployment}
\textbf{离线/校准阶段的完整算法。}
\begin{enumerate}
\item 使用已有策略训练/固定源条件 prior；保留多个源行为模式与 source dynamics 信息。
\item 在声明候选集内选择 $B$ 次校准交互，划分拟合与独立验证；拟合目标 context 和 ensemble 不确定性。
\item 检查动作接口与模式支持；固定接口，标记验证不足区域。
\item 采样 $q,g,m,U_s,\delta$，通过 $T_\varphi$ 得到 $U_t$；在源与 calibrated target model 中 rollout。
\item 计算目标任务、实际响应、模式与风险成本，更新 transport 与乘子；检查改变后的访问支持。
\item 在模拟验证机制与保留的校准响应上选择超参数；锁定方法后进行真实 held-out 评估。
\end{enumerate}

\textbf{在线阶段。} 每个技能边界读取实际 $q$，从 $p_s(\cdot\mid q,g)$ 采样，调用 transport，再通过 action decoding 执行固定短时技能；新的真实反馈用于下一次条件化。低层保持闭环。没有每次在线求解轨迹优化的必要。

若不确定性超过预设验证范围，则调用已有 hold/retreat/recovery 接口或停止。该事件单列为覆盖不足，并在未完成任务时计为失败。fallback 不能提供只对本文开放的额外任务监督。

\textbf{最小可行配置。} 单一共享 task-space 接口，4--8 个模拟机制，少量低维 context 参数，source-conditioned residual MLP，短时响应 surrogate，固定模式标签与统一 Lagrange 更新。验证这套配置后再增加 flow、多身体和视觉输入。这样可以先检验 insight，而不同时开发新生成模型、全新 RL learner 和大型 world model。

\textbf{与既有 work 的关系。}
\begin{center}
\small
\begin{tabularx}{\linewidth}{L{1.9cm}Y}
\toprule
Work & 在本文中的建模 insight\\\midrule
2GO & 任务阶段与活跃约束改变响应误差的几何和 transport 优先级。\\
MGA & 以实际 realization 构造 transport 目标，而非仅比较名义命令。\\
MDOC & 用可执行支持与已知硬约束组织探针和目标行为。\\
MD-COAS & 用 Lagrange 结构协调软语义目标、风险约束与预算。\\
FedGuide & 共享源行为 prior，同时允许目标机制下的局部适应，保留有效模式。\\
\bottomrule
\end{tabularx}
\end{center}

\section{Supporting Analysis}
\label{sec:theory}
本节给出辨识、响应误差和 transport 成本之间的条件链；这些条件本身需要验证，不能由实验成功倒推出成立。

\subsection{Task-Relevant Identifiability}
在固定局部接触模式中，设
\begin{equation}
 y=A_{\mathcal B}\theta+\xi,\qquad b_g=B_g\theta.
\end{equation}
\begin{proposition}
任务响应 $b_g$ 在无噪声校准中可辨识，当且仅当
\begin{equation}
 \ker A_{\mathcal B}\subseteq\ker B_g.
\end{equation}
\end{proposition}
它说明校准应覆盖影响 transport 的机制方向，不一定需要完整物理参数。核空间不满足时不能用伪逆 trace 虚构零不确定性。固定设计 Gaussian 分析可形成响应置信区域；适应性采样需要相应置信工具。该命题属于标准线性辨识，不单独作为原创贡献。

\subsection{Probe Coverage, Identification and Hybrid Events}
\textbf{不可辨识的边界。} 若两个机制对全部允许探针产生完全相同的观测 law，而其任务所需响应或可成功动作集合不同，任何仅使用这些探针的算法都无法保证在两者上正确适应。特别是只在接触时起作用的未知参数，不能仅凭未接触的观测辨识。自由空间辨识可以外推接触的前提是正确的结构模型与对相关参数的充分激励；未接触本身不意味着所有 gain/delay 参数均不可辨识。

\textbf{有限结构族的充分条件。} 对有限候选机制 $\Theta_F$，令固定探针集合的真实/模型响应 signature 为 $z(\theta),\hat z(\theta)$，且真实机制在候选集中。假设观测噪声范数不超过 $\eta$，所有候选的 probe-model 偏差不超过 $\epsilon_p$。若
\begin{equation}
 \min_{\theta\ne\theta'}\|\hat z(\theta)-\hat z(\theta')\|>2(\eta+\epsilon_p),
\end{equation}
则 nearest-signature 选择恢复真实候选：真实候选到观测的距离不超过 $\eta+\epsilon_p$，其余候选由 triangle inequality 严格更远。连续参数的离散化还需加入覆盖误差；当前 bounded profile search 不因此自动获得全局优化或统计置信保证。此结论解释一个受限结构模型的可验证机制，不声称辨识任意接触动力学。

\textbf{第一层机制验收。} CPU 首先固定已知身体/工具模型与 nominal mass，只辨识有界 servo gain 和整数 delay；公开 native joint friction 保留在身体模型中。从名义 Panda 学习并冻结事件条件 Gaussian reference prior，用计数的自由空间/加载探针拟合 twin，独立探针作拒绝检查，再在该 twin 上拟合受限 affine force transport。逆增益、名义响应 transport 与错误 context 对照用于区分普通参数补偿和实际响应优化。该层不替代 neural response ensemble、未知摩擦/顺应性、随机 transport 必要性或移动身体的正式验证。

\textbf{从校准到任务的缺口。} Probe 上的小误差没有直接推出下节的 uniform joint-kernel error。必须另声明结构模型正确性，或在 transport 访问支持上验证误差与覆盖；支持外概率与模型偏差不能省略。独立验证阈值是经验拒绝规则，不能称为概率安全证书。

\begin{proposition}[结构族内的校准到响应桥接]
令 $d_\Theta$ 为已声明机制族上的距离（可含离散 delay），真实机制为 $\theta_\star$。假设探针模型 signature 在该族上满足
\[
 \|\hat z(\theta)-\hat z(\theta')\|\ge\kappa d_\Theta(\theta,\theta'),\qquad \kappa>0.
\]
若观测噪声不超过 $\eta$，真实机制处 probe-model 偏差不超过 $\epsilon_p$，所选估计的实际拟合残差为 $r=\|\hat z(\hat\theta)-z_{\rm obs}\|$，则
\[
 d_\Theta(\hat\theta,\theta_\star)\le(r+\eta+\epsilon_p)/\kappa.
\]
进一步假设在共同访问支持及所有允许 transport 上，联合模型 kernel 对机制统一为 $L_\Theta$-Lipschitz，真实机制处的统一 kernel 偏差不超过 $\epsilon_K$，则下一节的误差可取
\[
 \varepsilon_{\mathcal B}\le\epsilon_K+\frac{L_\Theta}{\kappa}(r+\eta+\epsilon_p).
\]
\end{proposition}
\begin{proof}
第一式由 inverse separation 和经观测的 triangle inequality 得到；第二式在真实 kernel 与估计模型 kernel 之间插入真实机制处的模型 kernel，再应用 triangle inequality。该残差形式不要求优化器达到全局最优。
\end{proof}
这补齐了一个\emph{条件性}链条：结构族与任务支持正确、signature 可分、参数到联合闭环响应稳定，才可从探针外推到 transport。当前数值 profile 和独立验证没有认证 $\kappa,L_\Theta,\epsilon_K$ 的全局界。跨接触事件时统一 Lipschitz 条件可能不成立，应改用局部模式内界与下述事件 margin，不能把连续 gain 的小误差直接称为全任务保证。

\textbf{接触事件。} 设真实/预测轨迹有共同 coupling，轨迹距离的期望不超过 $e$。对一个成功/失败事件，若真实轨迹到事件边界的 margin 为 $M$，则任意 $r>0$ 下
\begin{equation}
 |\Pr(A)-\Pr(\hat A)|\le e/r+\Pr(M\le r).
\end{equation}
这是事件不一致只可能发生于距离至少 $r$ 或 margin 不超过 $r$，再应用 Markov inequality 的结果。孔口与力阈值附近的 margin 项可能很大；连续响应均值误差不能直接替代事件误差。接触建立、卡滞、恢复和超力需分别报告。

\subsection{Calibration Error and the Transport Objective}
分析给定初态的短时 transport 目标。令增广状态 $x$ 包含目标状态/历史、controller carry、事件 phase 与 tracking integral、源参考、源/目标技能指令和模式；真实与预测 rollout 使用同一个源参考及 transport sample。源参考仅是分析与训练成本变量，不向在线控制输入未观测目标未来。假设在同一个共同可行参数集合 $\Phi$ 内，所有 $\varphi$ 的联合增广闭环 kernel 都统一满足
\begin{equation}
 \sup_{\varphi\in\Phi,\,x\in\mathcal X_\Phi}
 W_1(\mathbb K_t^\varphi(x),\widehat{\mathbb K}_t^\varphi(x))
 \le\varepsilon_{\mathcal B},
\end{equation}
输出包含所用任务响应；支持覆盖与联合 kernel 的 $L$-Lipschitz 常数也在 $\Phi$ 内统一成立。仅有目标 marginal 接近不足以控制 source-coupled cost。相同增广初态下分布误差满足
\begin{equation}
 d_{k+1}\le Ld_k+\varepsilon_{\mathcal B},\qquad
 d_k\le\varepsilon_{\mathcal B}\sum_{j=0}^{k-1}L^j.
 \label{eq:propagation}
\end{equation}
对于由期望任务成本和 source-coupled trajectory cost 组成的目标，若相关成本具有统一 Lipschitz 常数，可得到
\begin{equation}
 |F_t(\varphi)-\widehat F_t(\varphi)|\le\varepsilon_F.
\end{equation}
在上述统一条件和共同可行集合中，若模型短时目标训练为 $\varepsilon_{\rm opt}$ 次优，则
\begin{equation}
 F_t(\hat\varphi)-\inf_{\varphi\in\Phi}F_t(\varphi)
 \le2\varepsilon_F+\varepsilon_{\rm opt}.
 \label{eq:objectivebound}
\end{equation}
这是所定义 transport objective 的辅助保证，包含 source anchoring 的偏置；不是对完整目标 task-only optimal policy 的 regret。CVaR、罕见事件和 UOT 增强目标需另外分析，不能直接由此推出真实安全或通用 OT 保证。

\textbf{改进何时可保证。} 对同一可行集合内的基线 $\varphi_0$，若模型成本改进 $\widehat F_t(\varphi_0)-\widehat F_t(\hat\varphi)>2\varepsilon_F$，则真实成本严格改进，直接由两端的统一成本误差界相减得到。若该差值不足，则上述分析不保证优于 direct transfer 或 system identification；随机性、flow 或 OT 名称本身也不产生改进保证。

\subsection{Feasibility Before Response Alignment}
\label{sec:finite-feasibility}
\textbf{实验要求补上的逻辑层。} P1 的 matched 条件中，校准模型几乎逐轨迹
复现了失败，但所选有限映射仍只有 $3/30$ 成功。因此，校准误差、控制/映射
表达误差和目标选择偏置必须分开。小的响应预测误差可以准确预测一个坏策略；
它既不保证候选类存在可行策略，也不保证 source anchoring 优于 task-only。
以下结论条件于固定接口、已声明模型与候选类，作为方法设计与审计工具。

\textbf{可行支持与拒绝。} 设 $\Phi_F$ 是冻结的有限映射族，
$S_m\subseteq\Phi_F$ 是模式 $m$ 下满足所声明任务、约束及模式对应的映射集合。
必须区分真实的 $S_m$ 与仅在有限模型训练/验证 reset 上观察到的
$\widehat S_m$；后者不是真实安全集合的证书。
\begin{proposition}[共同映射、模式条件映射与重加权的支持边界]
对于具有正源概率的模式集合 $\mathcal M_s$，一个共同映射能保留全部模式
当且仅当 $\bigcap_{m\in\mathcal M_s}S_m\ne\varnothing$。
允许逐模式选择同一候选族中的映射时，全部模式可执行当且仅当每个
$S_m\ne\varnothing$。仅对固定映射的源模式重加权不能使其不可执行模式
变为可执行；若全部模式均无支持，则输出停止/拒绝质量，而不是归一化零权重。
\end{proposition}
\begin{proof}
共同映射必须同时属于每个集合。模式条件映射从每个非空集合选择一个元素，
反向必要性直接成立。重加权只改变固定响应支持上的概率，不产生新的命令或
闭环响应支持，故不能补回缺失的可执行模式。
\end{proof}
这是\emph{候选类内}的结论；搜索失败不证明机器人全局不可行。
P1.7 的 task-only 也允许相同逐模式选择，因此增加条件化容量本身不能
作为 response anchoring 的独立优势。有限 reset 上的全成功过滤仅是经验
拒绝规则；所有拒绝与未执行任务保留在总体分母中。

\begin{proposition}[不变执行前缀的失败上界]
设一族 transport 在到达可观测阶段边界 $\sigma$ 前，采用相同指令、
相同控制器及反馈规则。将初态、源样本和物理随机流共同耦合，则它们在
$\sigma$ 前具有相同的增广闭环轨迹。若共同前缀在到达 $\sigma$ 前
以概率 $p_{\rm pre}$ 触发不可撤销终止失败，则该族任一策略的整体成功率
至多为 $1-p_{\rm pre}$。
\end{proposition}
\begin{proof}
相同初态和随机流下，按闭环递推归纳，相同反馈指令给出相同状态、
controller carry 与事件。故前缀失败事件对全族相同；整体成功事件必须
包含于其补集。对该包含关系取概率即得。
\end{proof}
这里的 $p_{\rm pre}$ 是声明初态分布下的概率，有限模型样本只给出其经验
描述。P1.6 模型训练中出现阶段 $1$ 的回撤超力，而阶段映射子族在阶段 $2/3$
才激活，说明应检查恢复前缀的表达支持。该结论不授权放宽力限或隐去失败；
新的恢复阶段参数族须单列版本，并对所有同类控制对照开放。
另做一个不回流搜索的配对模型诊断：同一已观察的失败 reset 中，六个只改
终段的候选在 $1.04\,\mathrm s$ 全部超力，失败前位置与指令逐元素相同；
加入回撤阻抗的 $18$ 个声明候选均成功。该单 reset 模型结果检验前缀
不变条件及表达缺口，不证明跨 reset、实际目标或总体的因果收益。

\textbf{保留模式质量与停止的代价。} 设 $S$ 是给定可行性规则允许保留
的源模式集合，$p_s(S)>0$。若目标 law $q$ 只在 $S$ 上分配质量，则
\[
 \operatorname{TV}(q,p_s)\ge p_s(S^c),\qquad
 \min_{q:q(S)=1}\operatorname{KL}(q\Vert p_s)=-\log p_s(S).
\]
第二个最小值由 $q=p_s(\cdot\mid S)$ 达到：对任意支持于 $S$ 的 $q$，
$\operatorname{KL}(q\Vert p_s)=
\operatorname{KL}(q\Vert p_s(\cdot\mid S))-\log p_s(S)$。
第一个式子直接在集合 $S^c$ 上使用 TV 定义。若保留原源采样而明确对
$S^c$ 的质量停止，则该停止策略的成功率至多为 $p_s(S)$。
这些是支持限制的基本恒等式，不作为新的 OT 理论；实验中仅有模型估计
$\widehat S$，不能把它当成真实全局可行集合。特别是未通过有限训练规则
不等于该模式成功概率为零；上述停止上界描述算法选择，而不是机器人能力上界。

因此，部分能力条件下不能同时要求源 law 不变、全部模式对应和高任务成功。
P1.7 的 reweight 对照单列检验这种取舍，不能因原始 source draw 的对应下降
就判为错误。若共同优化模式 law 与映射，应把 $(q,\phi)$ 作为联合候选；
下述固定 law/同集合的任务预算保证不能自动转移到不同重加权策略。

\textbf{用任务非劣预算组织响应目标。} 固定一个非空的共同候选集合
$\mathcal A$，令真实任务与响应成本分别为 $C,R$，模型成本为
$\widehat C,\widehat R$。先求 $\widehat C_*=\min_{\phi\in\mathcal A}
\widehat C(\phi)$，再选择
\begin{equation}
 \widehat\phi_\delta\in\arg\min_{\phi\in\mathcal A}\widehat R(\phi)
 \quad\text{s.t.}\quad
 \widehat C(\phi)\le\widehat C_*+\delta.
 \label{eq:task-band-transport}
\end{equation}
任务约束、mode 条件、候选容量和模拟缓冲对完整方法与 task-only 相同。
$\delta$ 在封存目标评估前固定；它公开允许的任务成本代价，而不是通过
反复调节响应权重寻找最终测试优势。

\begin{proposition}[同一有限集合内的任务预算与响应比较]
若在 $\mathcal A$ 上 $|C-\widehat C|\le\epsilon_C$、
$|R-\widehat R|\le\epsilon_R$ 统一成立，且 Eq.~\eqref{eq:task-band-transport}
求解误差至多为 $\epsilon_{\rm opt}$（约束仍被满足），记
$C_*=\min_{\phi\in\mathcal A}C(\phi)$，则
\begin{equation}
 C(\widehat\phi_\delta)\le C_*+\delta+2\epsilon_C.
\end{equation}
若 $\delta\ge2\epsilon_C$，还有
\begin{equation}
 R(\widehat\phi_\delta)\le
 \min_{\substack{\phi\in\mathcal A\\C(\phi)\le C_*+\delta-2\epsilon_C}}
 R(\phi)+2\epsilon_R+\epsilon_{\rm opt}.
\end{equation}
\end{proposition}
\begin{proof}
以真实 task minimizer 代入模型最小值，得到
$C(\widehat\phi_\delta)\le\widehat C_*+\delta+\epsilon_C
\le C_*+\delta+2\epsilon_C$。任何满足第二式真实成本约束的候选均满足
$\widehat C(\phi)\le C_*+\delta-\epsilon_C
\le\widehat C_*+\delta$，因而属于模型预算带；再对响应成本两端应用统一
误差界和优化误差即得。
\end{proof}
该命题只比较同一候选集合。本文开发实验中的 $\delta=0.02$ 是预先固定
的经验选择规则，并未认证统一 $\epsilon_C,\epsilon_R$，所以不能把实测
小误差直接代入，宣称总体非劣或真实安全。

\textbf{模型偏差与 reset 采样误差分开。} 上式的 $\epsilon_C$ 若将有限
训练均值当作 $\widehat C$，必须同时覆盖模型偏差与 Monte Carlo 误差。
对事先固定的 $M$ 个候选和校准后独立的 $n$ 个 reset clusters，
若每 cluster 的模式平均成本属于 $[0,1]$，则以概率至少 $1-\alpha$，可以取
\[
 \epsilon_C=\epsilon_{\rm model}
       +\sqrt{\frac{\log(2M/\alpha)}{2n}}.
\]
此处 $\epsilon_{\rm model}$ 仍需外部论证；有限 occupancy 的实测误差
不是它的统一证书。即使同一 twin 精确预测实际轨迹，四个训练 reset 上
的最优候选也可能在新模型验证 reset 失败。当前 $n=4,M=24$ 的保守界
远大于 $\delta=0.02$，所以预算带只是经验设计，不能宣称已认证非劣。
在 v2 的一个新 law reset 上，full 比 task-only 多一次失败；对应重复的
实际样本均值任务成本差为 $0.06788$，尽管响应成本差为 $-0.00309$。
这不是条件命题的反例，而是缺少统一误差与充分 reset 覆盖的实测提醒。
对有限样本的精确枚举最优，也不等于对期望目标的 $\epsilon_{\rm opt}=0$。

\textbf{物理后端误差单独计入。} 设 $C_{\rm CPU}$ 为冻结目标参考环境的
期望成本，$C_{\rm twin,CPU}$ 为校准参数在同一 CPU 物理下的成本，
$C_{\rm twin,GPU}$ 为这些参数在 GPU 物理下的成本。若分别存在同一有限
候选集合上的统一界，则由三角不等式，GPU 训练均值所需的统一误差可取
\[
 \epsilon_C = \epsilon_{\rm calibration}
                  +\epsilon_{\rm engine}
                  +\epsilon_{\rm sampling},\qquad
 \epsilon_{\rm engine}\ge
 \sup_{\phi\in\mathcal A}
 |C_{\rm twin,CPU}(\phi)-C_{\rm twin,GPU}(\phi)|.
\]
这里各项是需要论证的上界，而非由少数配对轨迹自动取得的证书。
坐标系、力矩平移、状态 carry 和随机流一致只是实现层前提；接触解算差异
仍可能通过阶段切换、首次保持终止或不可恢复失败放大。双精度本身也不保证
长时闭环等价。因此 GPU 模型搜索、独立 CPU 模型验证与冻结 CPU 目标评估
应分别记账。后两者的有限验证不能被解释为对所有候选与未来 reset 的统一界。

\textbf{表达误差不能被模型精度消除。} 对更宽的参考策略类
$\Psi\supseteq\mathcal A$，定义
$\Delta_{\mathcal A}=\inf_{\phi\in\mathcal A}C(\phi)-
\inf_{\psi\in\Psi}C(\psi)\ge0$。若原标量化目标
$\widehat C+\lambda\widehat R$（$\lambda\ge0$）的优化误差为 $\epsilon_{\rm opt}$，
$R\ge0$，则以同集合真实 task minimizer $\phi_C$ 为比较点有
\begin{equation}
 C(\widehat\phi)-\inf_{\psi\in\Psi}C(\psi)
 \le\Delta_{\mathcal A}+2\epsilon_C+2\lambda\epsilon_R
   +\lambda R(\phi_C)+\epsilon_{\rm opt}.
\end{equation}
证明由模型标量化最优性、成本两端误差以及
$-\lambda R(\widehat\phi)\le0$ 直接得到。即使校准和优化完全准确，
候选表达误差与 anchoring 偏置仍可能存在；上一节的
$2\varepsilon_F+\varepsilon_{\rm opt}$ 并未消除这些项。

\textbf{接触事件裕度与策略质量分开。} 上节事件界控制真实与预测的
事件差异，不保证事件本身成功。应同时报告任务失败概率、逐约束轨迹裕度、
以及预测事件误差。P1 的三个 matched 成功在最后保持窗口中距法向力下界
仅约 $0.0008$--$0.0144\,\mathrm N$，因此均值误差、最终力或模式标签
都不足以表示执行鲁棒性。事件距离必须包含接触 witness、计数器与控制 carry；
末端位置或平均力的 RMSE 不能替代它。
同时，控制器在首次满足保持条件时终止，会使最后保持窗口的起点接近阈值。
因此上述窗口裕度只是轨迹描述，不能直接当作闭环扰动鲁棒半径；验证后者
还需要控制扰动、完整停止时刻与必要的终止后承载诊断。

\textbf{验证拒绝属于算法成本。} 令 $a_{j,m}\in\{0,1\}$ 为目标评估前冻结
的模式接受标记，$C_{j,m}\in[0,1]$ 为该映射执行时的真实任务成本。
以拒绝成本 $1$ 计入全部请求时，实际算法成本是
\[
 C_j^{\rm alg}=\sum_m p_m\{a_{j,m}C_{j,m}+1-a_{j,m}\}.
\]
因此同一源 law 下，完整方法 $F$ 与 task-only $T$ 满足
\[
 C_F^{\rm alg}-C_T^{\rm alg}
 \le \sum_{m:a_{F,m}=a_{T,m}=1}p_m(C_{F,m}-C_{T,m})
       +\sum_{m:a_{F,m}=0,\,a_{T,m}=1}p_m.
\]
共同拒绝模式贡献为零；仅 task-only 被拒绝的项非正；仅完整方法被拒绝
的项为 $p_m(1-C_{T,m})\le p_m$，即得该式。训练预算带即使控制了
映射本身的成本，也不能忽略两方法的验证拒绝差异。若同时采用不同的
模式权重 $p_F,p_T$，按 $p_F$ 应用上述比较后还需加上
$\operatorname{TV}(p_F,p_T)$，因为停止后的成本仍属于 $[0,1]$。
这区分了被执行映射的质量、接受集合和模式质量分配，防止只在成功部署
子集上宣称总体非劣。若两方法具有相同冻结映射、接受标记与源 law，则
其算法成本由结构直接相等；这不提供独立任务收益。

\textbf{边缘一致不等于源条件对应。} 令 $M\sim\mathrm{Bernoulli}(1/2)$。
响应 $Y=M$ 与 $\widetilde Y=1-M$ 的边缘 law 完全相同，但
$\Pr(Y\ne M)=0$、$\Pr(\widetilde Y\ne M)=1$。
两种策略可以具有相同成功率与目标模式频率，却具有相反的联合对应。
固定打乱 coupling 的实验验证这一度量边界；它由构造产生，不能单独证明
本文方法优于使用同一 source prior 的 task-only 控制器。

\textbf{有限类选择持平的解释。} 对冻结成本 $(C_i,R_i)$，候选 $i$
最小化 $C+\lambda R$ 当且仅当
$(C_i-C_j)+\lambda(R_i-R_j)\le0$ 对全部 $j$ 成立。
这些半直线的交集给出可直接计算的 $\lambda$ 区间，包含空区间和无界区间。
同一个映射在共同随机流下产生相同实际 law；给它不同方法名称不会产生
独立收益。该代数诊断说明何时 anchoring 不改变选择，既不是统计证书，也
不支持通过试遍 $\lambda$ 后挑选测试集胜者。

\textbf{相同控制类与额外随机性的边界。} 本文 full 与强 task-only 使用
相同参数化与信息，因此可表达相同闭环 law；source anchoring 改变的是
选择准则，不扩大相对该强对照的可执行集合。相对 reweight-only 的支持
改变不能转称相对同类控制器的表达优势。

固定一个源样本、可观测条件和有限合法动作集合，令 $c(a)$ 为已对物理
随机性取期望的标量成本。任意条件随机 kernel $T$ 满足
$\sum_a T(a)c(a)\ge\min_a c(a)$，确定性最小动作即可达到等号。
因此在可逐条件分解的期望目标中，额外 transport noise 没有自动收益；
随机源 prior 经确定性映射仍可多模态。分布耦合约束、条件覆盖需求或
不同的鲁棒博弈可能改变这一结论，必须单独分析。当前 extra-noise 对照
只是固定小扰动，并非训练过的同容量随机 kernel，不能据其单次成功数
差异证明随机 transport 必要性。


\subsection{Finite Frozen-Policy Audit}
对预先冻结的有限策略集合，可以采用较弱但可执行的经验审计，而不宣称验证了上述无限策略类的统一 kernel 假设。设策略数为 $M$，校准后独立抽取 $n$ 个 reset clusters，每个 cluster 内平均所有预声明模式；同一 cluster 的实际与模型成本使用匹配随机流。成本函数在评估前固定并有界于 $[0,1]$，记观测为 $a_{ji},p_{ji}$，其均值为 $\bar a_i,\bar p_i$。令
\[
 \beta_n=\sqrt{\frac{\log(4M/\alpha)}{2n}},\qquad
 \epsilon_n=\min\!\left\{1,\max_i\frac1n\sum_j|a_{ji}-p_{ji}|+
 \sqrt{\frac{\log(2M/\alpha)}{2n}}\right\}.
\]
\begin{proposition}[有限冻结策略的成本改进审计]
若 clusters 独立同分布且与策略选择数据独立，则以至少 $1-\alpha$ 的概率，对所有冻结策略 $i$，
\[
 F_t(\varphi_0)-F_t(\varphi_i)
 \ge\bar p_0-\bar p_i-2\epsilon_n-2\beta_n.
\]
\end{proposition}
\begin{proof}
对 $M$ 个模型均值使用双侧 Hoeffding 界与 union bound，对 $M$ 个有界绝对配对误差使用单侧 Hoeffding 界与 union bound，各分配 $\alpha/2$。由 $|F_t(\varphi_i)-\widehat F_t(\varphi_i)|\le\mathbb E|a_i-p_i|$，再对基线和策略两端应用 triangle inequality 即得。
\end{proof}
该结论条件于一次校准、声明的 reset 分布、成本和有限策略；多次校准需另外分层或分配置信预算。小样本时右端可能为负，必须返回未认证。该界不能覆盖未冻结的 neural policy、支持外接触、任意实机噪声或任务成功指示函数之外未记录的风险，也不保证相对采用不同目标的控制器必胜。

\paragraph{直接观测的成本改进。}
若只验证冻结策略的实际效果而不认证模型，可以对独立 reset 的配对成本差 $d_{ji}=a_{j0}-a_{ji}\in[-1,1]$ 使用更直接的单侧界。对 $M-1$ 个预声明比较，以至少 $1-\alpha$ 的概率同时有
\[
 F_t(\varphi_0)-F_t(\varphi_i)
 \ge \bar d_i-\sqrt{\frac{2\log((M-1)/\alpha)}{n}}.
\]
证明直接对长度为 $2$ 的有界配对差应用 Hoeffding 与 union bound。该式只认证固定成本的实际改进，不能代替校准可辨识性、response kernel 精度或跨校准统计；样本量与比较集合须在这些新 reset 执行前固定，禁止 optional stopping。

\section{Experiments}

\paragraph{后续六项机制验收。}
旧的宽开口 CPU fixture 与下面的匹配装配 pilot 分开保存。新 \texttt{matched\_keyway} 在 $20\times12\,\mathrm{mm}$ 主 peg 和 $+x$ 凸键周围构造专用键槽，单边间隙开发起点为 $0.8\,\mathrm{mm}$；错误键位必须真实碰壁，reset 不允许穿透。对齐使用已知几何投影的 peg/key 全长包络及正余量，而非仅检查 TCP 圆形误差。插入过渡中的间歇侧触与入座后的轴向支撑丢失分开判定，所有方法共用该事件语义和原力限/时限。新的 wrench 选择作用于工具的接触力，并以接触点到 TCP 的 moment arm 平移力矩，均在 world 坐标中表达；旧 aggregate 成绩不因此自动成为完整 F/T 证据。

验收依次为：(i) 物理与时间步/源模式资格；(ii) 同一校准模型、同一候选池的控制搜索、task-only residual 与短时 MPC 对照；(iii) 接触摩擦/顺应性、模型偏差、读出噪声与任务 occupancy 预测检查；(iv) 实际采样的模式 law 和额外随机性消融；(v) 多次独立校准、身体$\times$机制交叉及 UR5e 留出；(vi) 主张/实现/证据矩阵和 Task 1/2/4 增量。读出噪声与闭环传感噪声是不同协议。前三项有限结构研究不替代 learned response 的验收。

Mirage 已研究考虑动力学后的任务空间跨臂迁移\cite{mirage}，FADA 采用冻结 planner 并以少量目标交互适配 inverse dynamics\cite{fada}。因此冻结 prior 和少量校准本身不足以支持新颖性主张；本工作的增量必须由任务相关接触响应、约束内分布迁移及同预算强对照建立。若同模型控制搜索解释全部收益，或额外噪声没有独立收益，应相应收窄贡献。详细 go/no-go 与实现索引见 \url{docs/development/cross-embodiment-theory-feasibility.md}。
\label{sec:experiments}
以下为正式评估协议，不填写预期优势。实验从真实形态的全 3D 机械臂开始，依次增加受限空间运动、受约束接触和装配恢复，再用 Task 4 将同一操作扩展到移动身体；不设置 2D 或平面点质量实验。Task 1--3 建立机制证据，Task 4 检验身体变化扩大后的适用范围。Preliminary 中的解析 gain-shift 例子只解释指令 law 与响应 law 的差别，不充当实验结果。

\paragraph{CPU 工程验证与正式评估分开。}
当前小型实现已接通 Panda 源先验、合法源状态投影、响应 ensemble、预算校准、随机残差 transport 与真实物理回放，包含固定臂导向/插入、平面移动底座以及 Go2+Z1 站立操作 adapter。接触力来自仿真约束；浮动身体额外保存足锚点、QP warm start 和支撑失败。CPU 的 JAX 整段执行与 Torch/C++ 数值接口验证用于建立可复现代码基线；小样本成功、保留的超力/超时失败以及未校准 ensemble 均不构成机制优势证据。探索性目标配置已用于调试，正式结果需另冻结测试 split。Task 4 继续复用 Task 3 的冻结 Panda prior；平面底座的高度/侧倾约束、Go2 的站立控制范围与 host QP 成本须明确披露。GPU、正式多校准种子/未知接触统计和人形接口尚未验收。实现与结果工件索引见 \url{docs/development/cross-embodiment-transport-plan.md}。

\paragraph{历史宽开口 CPU 机制协议。}
第一层固定臂机制研究使用显式命名的 research 配置，区别于下面的拟议工业工件 pilot：共同刚性 peg 为 $20\times12\times70\,\mathrm{mm}$，带附加 key；socket lip 高度为 $0.481\,\mathrm{m}$，目标有效入座 $16\,\mathrm{mm}$。位置误差 $<2\,\mathrm{mm}$ 保证深度 $>14\,\mathrm{mm}$，方向误差 $<2^\circ$，法向力 $5\pm2.5\,\mathrm N$，连续保持 $0.3\,\mathrm s$；子步合力峰值不超过 $25\,\mathrm N$、侧向合力不超过 $10\,\mathrm N$，含恢复的 timeout 为 $15\,\mathrm s$。初始 lip 偏置为 $10\,\mathrm{mm}$，另加均匀 $\pm2\,\mathrm{mm}$ 任务空间扰动、向上 $1$--$3\,\mathrm{mm}$ 偏置，以及 $0.12\pm0.01\,\mathrm{rad}$ yaw。mode0 直撤，mode1 侧撤；模式由实际接触与撤离轨迹分类，不能以请求标签代替。

所有目标对照共用 decoder、$4\,\mathrm{Nm}$ 角向 wrench 上限、$15\,\mathrm{mm/s}$ 接近速度和有界接近积分；真实 actuator/关节/速度/碰撞限制不变。5 次计费探针（3 fit、2 validation）每次最多 $2\,\mathrm s$，自由空间探针提供足够的平移/角向激励，先作 gain/delay profile，再局部联合拟合接触响应。所有搜索调用另报；参数搜索无全局证书。源 Gaussian reference prior 冻结后，阶段条件 affine transport 只调整最终插入阶段的法向力倍率和切向速度偏置；源恢复阶段不随最终插入力一起放大。优化仅在 calibrated/public twin 的 2 个训练初态、2 种模式上进行，固定 48 次 rollout；nominal-response 对照具有同样额度。

这套历史宽开口配置用于检验正确结构族内的机制链，不替代未知接触模型与工业夹具测试。其法向观测定义为约束接触载荷 $\sum_i|F_{i,z}|$，入座保持阶段为单向 floor 支撑；当时接触转矩未加 TCP moment-arm 变换，不能当作完整 wrist F/T 资格验证。上面的 matched-keyway 协议已改用工具 TCP world wrench，两套数字不可混合。更早的 CPU 配置约 $46\,\mathrm{mm}$ 深度、无姿态检查及较长 timeout 结果也保留为历史诊断。新结果须报告独立 reset-cluster 的成功率/区间、模式、姿态、子步力峰值，以及所有被验证拒绝的机制；独立验证阈值不因失败而放宽。

环境按操作语义命名：ObstacleTransfer、GuidedContactAlignment、KeyedInsertionRecovery、MobileKeyedInsertionRecovery。代码/config identifier 分别为 \texttt{obstacle\_transfer}、\texttt{guided\_contact\_alignment}、\texttt{keyed\_insertion\_recovery}、\texttt{mobile\_keyed\_insertion\_recovery}；Task 1--4 仅表示本节科学证据的排序。

\subsection{RQ1: Cross-Embodiment and Sim-to-Real Effectiveness}
\input{sections/current_evidence.tex}

\textbf{机器人与已知信息。} 源臂为 Panda（7 DoF）；xArm7（7 DoF）为工程桥接目标，复用已有 robot profile，但仍需验证新任务和控制 adapter。UR5e（6 DoF）已接入 CPU 模型、关节索引和控制接口；正式跨自由度结果仍须通过完整评估。使用公开 UR5e 响应预训练的工程对照属于已见身体，不能并入下文未见身体主协议。主设置使用相同工件、夹具、工具和任务坐标系。目标 URDF、连杆/持件碰撞几何、关节边界和声明控制能力作为已知信息向全部方法开放；未知 controller/contact context 包括延迟、死区、有效增益/阻尼、持件惯量误差及接触摩擦/顺应性。已知形态差异不包装成需要干预发现的潜在机制。

\textbf{统一的新技能 adapter。} 四项任务固定使用 $U\in\R^{H\times13}$，每步为
\begin{equation}
 U_j=\big[v_j^\top,\ \omega_j^\top,\
 (\log k^{\rm trans}_j)^\top,\
 (\log k^{\rm rot}_j)^\top,\ F^d_{n,j}\big]^\top.
 \label{eq:experimentalprimitive}
\end{equation}
其中 $v,\omega\in\R^3$ 为任务空间线/角速度，六个 log-stiffness 分别控制三平移和三转动方向的对角刚度，$k$ 表示各轴刚度与固定参考值之比，$F_n^d$ 为当前任务接触面的法向力设定值。坐标系、单位、归一化、技能时长、动作边界及刚度范围预先固定；自由空间阶段将力维 mask 为零，接触阶段的法向由已知夹具几何和可用状态确定。各臂通过自己的运动学和低层控制器落地这些共同指令；所有方法使用同一目标 adapter。该 13D 接口需新实现，不声称与现有 10D 扫描 primitive 或其他同维接口语义兼容。若实机缺少某项控制能力，应在共同可执行子空间重新定义并重跑全部对照。

\begin{enumerate}[label=\textbf{Task \arabic*.},leftmargin=2.1em]
\item \textbf{ObstacleTransfer（持件绕障转移）。}
机械臂固定夹持带方向性的 L 形工件，穿过包含横梁、侧板和姿态受限开口的三维工位，到达指定装配前 $SE(3)$ 位姿。目标同时要求位置与方向；检查全臂、夹爪和持件碰撞、关节边界及速度/力矩约束。源 prior 保留上绕与侧绕两类路径和相应腕姿态，匹配初态/目标下均有成功源轨迹。已知形态改变影响可达与通行集合；未知执行延迟、惯量或饱和改变制动、转向和跟踪响应。短探针为多个非奇异配置下的平移/转动脉冲、变速跟踪、停止与反向；检验目标是否需要移动路径支持并改变时间结构，而非仅调整单步增益。

\item \textbf{GuidedContactAlignment（导向接触对齐）。}
将带方向性的工件导入倾斜、具有三维侧向导面的定位夹具，借助接触滑移与旋转完成指定装配前位姿。主设置包含两种成功源接触序列：沿不同导面建立接触并滑入同一定位区域；运动同时涉及高度和姿态，不退化为桌面平面推物。保持工件/夹具几何，分别改变目标臂控制响应和接触摩擦/顺应性，再测试组合变化。探针是在有限压靠范围内的短时法向加载/释放、切向微移与小角度调整，记录力、滑移、接触丢失和工件位姿进展。核心设置使自由空间跟踪接近、接触响应不同，检查是否必须调整压靠方向、速度、刚度及接触序列的概率。

\item \textbf{KeyedInsertionRecovery（带键插入与恢复）。}
将固定夹持的非轴对称 connector 对准并插入对应 socket；从预声明的横向/角度偏移启动，并在部分试次中施加随机指派、幅度受限的 fixture 偏移或外扰。源 prior 包含两类有效分支，例如小角度搜索后推进与受限回撤--重对准--再推进，具体标签依据执行中的接触/位姿事件固定。区分机器人控制延迟/死区、工具柔顺性与 fixture 摩擦变化，不能把所有变化统称 embodiment。探针限于孔口附近的横向/旋转微扰和短推进/回撤；不得免费使用完整目标成功装配演示。除了最终插入，重点检验 jam-to-recovery 概率、重试后的完成及峰值侧向力，避免用低平均力掩盖停滞。

\item \textbf{MobileKeyedInsertionRecovery（移动身体带键插入与恢复）。}
复用 Task 3 的 connector/socket、工具、插入深度、间隙、力限与恢复模式，把同一工位操作从固定 Panda 扩展到移动身体。工位保持固定，观察接触反作用力、身体姿态变化与支撑响应如何改变 TCP 跟踪、卡滞和恢复概率。13D 指令不增加底盘速度、足步或步态标签；已知全身控制器落地相同末端/阻抗/法向力参考，transport 调整任务技能的时序、柔顺性与恢复分布。目标是在相同接口和全身控制器下验证响应校准是否仍有增量收益。
\end{enumerate}

\textbf{Pilot 起始数值：正式测试前冻结。} 下表是可调整的开发起点，不是已验证规格。只在声明的训练/开发机制上检查接触稳定性、可达性和时间步收敛；在访问 held-out 机制与最终评估前冻结几何、初态分布、阈值和 probe 时长。位姿误差分别使用欧氏位置误差与 $SO(3)$ 测地角；成功还须满足全程碰撞/关节约束和相应力限。表内力限是仿真实验判据，不能作为实机安全阈值；真实平台的执行界限需另行验证、预注册。
\small
\begin{longtable}{L{1.0cm}L{5.5cm}L{5.5cm}}
\toprule
任务 & 场景与初态起点 & 成功判据与停止条件\\\midrule\endfirsthead
\toprule 任务 & 场景与初态起点 & 成功判据与停止条件\\\midrule\endhead
Task 1 & 工位包围盒 $0.8\times0.6\times0.6\,\mathrm{m}$；L 形件包围盒 $120\times80\times20\,\mathrm{mm}$；持件质量 $0.3$--$1.0\,\mathrm{kg}$。 & 目标位置误差 $\le5\,\mathrm{mm}$、方向误差 $\le3^\circ$，保持 $0.3\,\mathrm{s}$；$10\,\mathrm{s}$ timeout；全臂/持件碰撞或声明控制约束超限计失败。\\
Task 2 & 工件 $80\times50\times20\,\mathrm{mm}$；单边导向间隙 $2$--$4\,\mathrm{mm}$；初始横向偏移 $\pm10\,\mathrm{mm}$、角度偏移 $\pm10^\circ$，具体轴与抽样分布在 pilot 后冻结。 & 定位误差 $\le2\,\mathrm{mm}$、方向误差 $\le2^\circ$；接触合力 $\|F\|\le20\,\mathrm{N}$；$10\,\mathrm{s}$ timeout；超力、碰撞或停止未定位计失败。\\
Task 3 & 刚性 keyed 件 $30\times18\times12\,\mathrm{mm}$；入座深度 $15\,\mathrm{mm}$；单边间隙 $0.5$--$1.0\,\mathrm{mm}$；初始横向偏移 $\pm2\,\mathrm{mm}$、角度偏移 $\pm3^\circ$。 & 插入深度 $\ge14\,\mathrm{mm}$、方向误差 $\le2^\circ$，保持 $0.3\,\mathrm{s}$；合力 $\le25\,\mathrm{N}$、相对插入轴侧向力 $\le10\,\mathrm{N}$；$15\,\mathrm{s}$ timeout，含恢复时间。\\
Task 4 & 保持 Task 3 工件/夹具和任务相对初态；增加已知底盘/支撑姿态与预声明的身体扰动。站位从共同任务支持中生成；平台尺度和扰动上限由独立 pilot 冻结。 & 沿用 Task 3 的插入、力限和 $15\,\mathrm{s}$ 操作 timeout；另须满足平台专属稳定/滑移/关节约束。倾覆、跌倒、非预期身体碰撞或超限均计整体失败。\\
\bottomrule
\end{longtable}
\normalsize
机制 pilot 采用有效伺服增益倍率 $\{0.7,1.0,1.3\}$、指令延迟 $\{0,20,60\}\,\mathrm{ms}$；Task 2/3/4 的工件接触摩擦系数起点为 $\{0.15,0.35,0.6\}$，静/动摩擦约定与接触离散化在测试前固定。Task 4 的轮地/足地摩擦与身体跟踪等级另做稳定性 pilot，不直接复用工件摩擦范围或放大伺服偏差。先做单因素，再做预声明组合；各等级在身体之间交叉，禁止仅按机器人身份绑定某个机制等级。工位包围盒不表示所有点都可达，起终位姿从共同可达区域预生成，通道/导面几何仍需保留要测的形态差异。

允许的夹持、工件--导面及 connector--socket 接触由白名单声明；Task 4 另声明轮地/足地支撑接触。碰撞失败指其他非预期接触。力峰值/超限在物理子步监测，不能靠控制频率降采样隐藏尖峰。Task 3/4 的 jam 事件起点定义为连续 $0.5\,\mathrm{s}$ 轴向进展小于 $0.2\,\mathrm{mm}$ 且存在持续接触载荷；载荷阈值在 pilot 后冻结。恢复初态由可重复的受限接近产生，不通过穿透几何制造卡滞；初态不可恢复的比例另报。

\subsubsection{Task 4 Protocol: Same Operation, Mobile Body}
\label{sec:task4}
\textbf{源与目标身体阶梯。} 源始终是冻结的 Panda Task 3 prior，不增加目标成功示范。第一层先比较固定 xArm7 与相同臂/工具安装在轮式底盘上的仿真组合，以隔离移动支撑的作用；底盘质量、轮地接触、安装变换和全身 controller 均需新增并验证，不能假定已有组合资产。第二层为带臂四足，候选为 Go2+Z1 仿真组合；第三层为具有可用夹持器的 G1 单臂操作配置。单独机器人模型不等于已验证的组合系统，安装几何、质量/惯量、工具、可达姿态、足端接触和 controller 均需资格测试；模型不具备共同技能能力时报告不支持，不临时放宽任务阈值。至少完成轮式平台和一种腿式平台后再主张跨移动支撑类型；人形配置按可用接口增量接入。

\textbf{同任务与支持匹配。} 主协议从工位附近的配对任务状态开始，保持工件--socket 相对位姿、初始速度、接触阶段、工具和目标一致；各身体的关节、基座和支撑状态可以不同。源 virtual state 必须在 Panda 的已验证支持内存在，接触历史通过共同 warmup 规范建立，不能用 IK 复制目标的身体/接触记忆。Task 4 延续 Task 3 的搜索与回撤--重对准模式，不把步态类别当作新的任务模式。远距离到站或超出源初态支持的移动接近单列为扩展：共享相同已知导航/站位模块，报告到站及操作的完整失败分母与耗时；不能把更大可达范围自动解释为原源分布的成功迁移。

\textbf{执行与身体状态。} 四轮/双轮等底盘形式须在 pilot 前选定并冻结；第一层优先选易于验证稳定性的轮式底盘。轮式全身 decoder 协同分配基座/机械臂运动；腿式全身 decoder 协同分配躯干、机械臂及允许的支撑调整。先用 Task 2 的限力接触检查站立/支撑与力接口，再运行不改变 Task 3 成功判据的正式扩展。自由空间接近可允许预声明的小幅移动或换步，接触插入先在明确支撑阶段执行，基座/躯干仍保留真实动力学和力反馈。所有受控方法使用同一目标 WBC、支撑切换规则、站位模块与可行性投影；WBC 不替 transport 搜索插入/恢复序列，投影与饱和须记录。目标身体条件加入基座 twist、相对工位姿态、支撑接触/滑移、关节与执行余量和 controller history，响应 surrogate carry 同步扩展；这些字段不直接塞进冻结 Panda prior，源查询仍只接收其声明的语义/virtual-state 输入。身体专属稳定特征仅作条件与目标风险；共享 response matching 不惩罚为完成任务所需的合法身体运动。若现有 WBC 只支持末端位姿而不能实现声明的阻抗/力接口，应对该 body 统一降维并重跑全部方法，单列为接口变体。

\textbf{机制变化与交叉控制。} 比较固定/释放基座，以及移动身体内的名义/偏移执行响应；基座锁定仅作诊断，不混入移动身体主成功率。未知 context 分别改变有效基座/躯干跟踪增益与延迟、轮地/足地摩擦、载荷质量/质心误差，再与 Task 3 的 wrist/contact 机制组合。名义几何、质量模型、支撑规则及控制能力均作为已知信息，校准只估计实际偏离。主机制 cell 尽量匹配自由空间 TCP 跟踪，却保留接触时的身体顺应与恢复差异；检查普通任务空间动作校准是否已足够。预声明小幅躯干/基座外扰可用于检验鲁棒性，但扰动指派与上限对方法相同且不参与目标调参。

\textbf{校准探针与预算。} 保持总预算 $B\in\{0,10,30,100\}$，在同一总时长上分配 (i) 末端脉冲/制动，(ii) 轮式小幅基座平移/转向或腿式稳定支撑内的躯干姿态变化，(iii) 孔口限力压靠与小切向移动的耦合探针。第二类是声明的辨识干预，不新增到部署技能维度；所有方法具有相同探针接口。联合记录 TCP、wrench、基座姿态/twist、支撑接触和滑移；腿式不使用与其模型不符的轮式速度命令。接口校准、WBC 的目标域调整、支撑探测、失败、reset 和 held-out validation 都计入同一 $B$。预训练 WBC/响应模型的公开机制数据和成本单列，隐藏目标只提供预算内 probes 与封存评估；目标域重训 WBC 属于额外适应协议，不能混入主表。

\textbf{关键对照。} 在相同目标 WBC 上比较 direct prior、action calibration only、shared prior+context、deterministic/full transport，以及同一 calibrated model 上的预算匹配 MPC/MPPI。增加 arm-only calibration 与完整 coupled calibration，共享总预算和时长；再增加不输入身体/支撑历史、以及不建模基座--接触耦合的 response 消融。仅 TCP 的版本只从高层网络移除身体条件；公共 WBC 和约束检查继续使用同一身体反馈，避免同时削弱低层控制。基座锁定和已知真值响应只作诊断。UMI-on-Legs 提供操作策略与 WBC 分层的参考\cite{umilegs}；如适配其接口或 controller，须明确改编、相同监督与 WBC 条件，不能将使用额外任务示范的原系统并入无示范主表。

\textbf{指标与结论边界。} 主指标是满足 Task 3 力/插入判据且全程保持身体约束的整体成功率；同时报告 jam/recovery 概率、实际模式频率、峰值侧向力、基座/躯干姿态偏移、足/轮滑移、倾覆/跌倒、WBC tracking/projection/saturation 率与端到端延迟。平台稳定裕度使用适合其动力学和支撑阶段的定义并分别报告，不将固定臂的零基座运动当作目标。所有操作 reset 均进入分母，稳定但未插入仍失败；到站扩展另报整段成功率。分别报告移动 body 是否参与响应预训练，以及 frozen-transport context-only 与 calibrated-model few-shot adaptation，沿用 UR5e 的数据权限区分。若共享 WBC 或普通动作校准已解释全部收益，则 Task 4 只支持接口可迁移；只有同 WBC 下有限校准改善耦合响应预测并改善受约束任务/模式结果，才支持本文机制的跨移动身体扩展。

\textbf{身体$\times$控制器与跨 DoF。} Panda/xArm7 分别配名义与偏移控制器，构成同身体不同响应、不同身体相近响应、不同身体不同响应的交叉矩阵；接触参数变化单独标记。UR5e 不进入源 prior、响应模型或 transport 的目标真值预训练，两个协议分开报告：(i) \emph{frozen-transport generalization}，冻结 $\varphi$ 与共享模型权重，仅用预算内记录拟合 $c_t$；(ii) \emph{held-out-body few-shot adaptation}，在拟合 $\hat c_t$ 后，允许用校准 twin 和匹配的离线预算更新 $\varphi$。二者均可使用目标已知形态；不能把更新 $\varphi$ 后的结果称为冻结 transport 泛化或零交互迁移。几何 OOD、机制 OOD 及其组合分别报告。

\textbf{实物与范围。} 实机型号、控制权限、腕部 F/T 和工具接口待确定。先验证新 adapter，再按同一工位协议开展可实施任务；不能预先把 Nova5 或任何现有实机接口当作已接通。仅有一个真实目标臂时，结论限制为 sim-to-real 或 cross-arm sim-to-real。Task 4 已纳入仿真扩展协议；移动身体实机需另外确认 WBC、工具、力觉与支撑状态接口，不能用仿真结果替代实机验证。

\textbf{任务指标。} 四项任务均报告无碰撞/超限的整体成功率、完成时间、任务成本、停止/超时比例和推理延迟。Task 1 增加最终 $SE(3)$ 误差、最小全臂/持件 clearance、跟踪误差与饱和率；Task 2 增加定位误差、接触保持、滑移/停滞率及力峰值；Task 3 增加插入深度、卡滞率、条件恢复概率、重试次数与侧向力/力矩峰值；Task 4 另加身体稳定与移动--接触耦合指标。全部 reset 均进入整体成功率分母，恢复条件成功率只能作为补充。

\subsection{RQ2: Calibration Budget and Generalization}
\textbf{预算与数据隔离。} 预算曲线首先使用 $B\in\{0,5,10,30\}$ 次目标短 probe；是否增加到 $100$ 次由预声明开发曲线决定，正式测试前固定。每次尝试按一个固定时长上限和一次 reset 计数；接口调试所需目标执行、fit probes、独立 held-out 响应验证和失败尝试全部纳入总数，每次只计一次。预先固定 fit/validation 分配，验证样本不参与 context 拟合；不能通过验证失败后免费补采绕过 $B$。$B=0$ 不访问目标真值，使用声明的名义/随机化先验。最终任务评估在模型和超参数冻结后独立进行，不回流训练，并另报总执行时间与 reset 成本。

\textbf{真值目标与训练 twin 分开。} 隐藏真实参数的 target simulator 只生成预算内 probes，以及封存的最终任务/诊断评估。拟合 $\hat c_t$ 后，transport、策略和 surrogate 的训练只能查询由已知身体模型与 $\hat c_t$ 构造的 calibrated surrogate/twin；不得直接用真值目标 simulator 做 GPU rollout、RL 更新、梯度计算或超参数搜索。最终评估及可行性诊断的查询量单列统计，其轨迹、参数和标签均不回流训练。所有方法遵守同样访问权限。

\textbf{探针设计对照。} 在相同可行候选集合与总时长下，比较 random probes、已有 task-dependent design、全参数辨识设计和只有自由空间探针的校准；接触任务增加相同总预算下的自由空间/接触混合分配。上述设计对照固定 source prior、响应模型结构与 transport；不把 task-aware 探针设计本身作为新贡献。相同 buffer 对照中，每种方法获得完全相同的可用记录；原生主动辨识系统可自行选探针，但其全部目标尝试受同一 $B$ 限制。

\textbf{泛化与指标。} 留出初态/目标位姿、未见控制器参数组合、持件惯量及接触条件，分别评估已声明机制族内泛化与 outside-support 接触分支。报告独立短时响应误差、力/进展预测、模式概率校准、模型覆盖、held-out discrepancy、成功率--预算曲线和真实分钟；不能仅用全参数估计误差代表任务响应校准。所有新增校准模型 rollout 与离线策略更新均报告并计入共同模拟/优化预算，不能将目标真值模拟器当作免费训练环境。

\subsection{RQ3: Distribution Correspondence, Modes, and Feasibility}
\input{sections/current_law_evidence.tex}

\textbf{源模式与目标能力。} Task 1 使用三维绕行类别；Task 2/3 使用实际接触序列或恢复事件；Task 4 继承 Task 3 的操作模式，并单列身体支撑事件，不将步态变化充当操作模式覆盖。标签和响应特征在实验前固定，source prior 在同一 $q,g$ 下提供多模态成功行为。接触主设置要求检验两种源模式在目标上经调整后均可执行的情形；另设能力变化诊断，检查部分模式不能保持时的质量重新分配。已知几何约束与校准模型可供训练使用，目标真值可行性标签不得进入 transport、mode loss、探针选择或超参数选择。

\textbf{可行性 oracle 只作诊断。} 在全部训练冻结后，用目标真值机制和较高计算预算的参考控制/规划器独立搜索、执行并验证各模式。只有找到满足约束且成功的行为才标为 verified feasible；搜索失败标为未验证，不充当不可行证明。oracle 不提供目标训练演示、模式 mask 或 replay。若预期双模式 cell 未得到双模式验证，保留其整体任务结果并如实标注；不事后删除困难 cell。模式覆盖在已验证模式集合内计算，同时公开未验证比例及全部任务分母。

\textbf{分布指标。} 对匹配初态/任务和相同 reset 分布的独立闭环 rollout，报告 verified-feasible mode coverage、实际 mode frequencies、接触/恢复事件概率、约束违反质量，以及归一化任务响应特征上的 Wasserstein/Sinkhorn discrepancy。展示指令分布与实际响应分布两个层次，并按模式分别报告形变、时间及任务成本。真实系统使用经验估计和 bootstrap，不称完整 law ground truth；必要的模式改变未必减小全局 source--target 距离。不能用 t-SNE、flow 可逆性或训练 loss 代替模式验证。

\textbf{受控消融。} 比较完整 stochastic transport 与同容量、同目标的 deterministic transport；后者仍可通过随机源 prior 产生多峰输出，不能将其预设成单峰。使用完全相同的 $U_s,q,g$、reset 和扰动列表做 $\varepsilon$ 消融，另固定单个 $U_s$ 重复采样，测量额外噪声是否增加可执行的条件模式覆盖或改善任务结果；条件熵升高本身不是收益。允许实验得出额外噪声无必要的结论。再比较 target-task-only residual（$\lambda_s=\lambda_M=0$）、仅去掉响应匹配（$\lambda_s=0$）、reweighting-only、nominal feature alignment、普通多步响应模型、uniform geometry、无 mode conditioning 与无风险约束。每项固定 probe buffer、source 初始化、目标 adapter、传感器、离线模拟量与更新次数；conditional flow 仅在带来独立分布收益时才作为必要组件。

\input{sections/current_gpu_evidence.tex}

\subsection{Existing Baselines and Fair Protocols}
\textbf{受控主表。} 协议 A 不使用目标完整任务演示或目标在线 RL fine-tuning。所有受控方法从同一 source prior/训练数据出发，获得相同已知形态、传感器及目标 adapter；允许在各自校准模型中做预算匹配的离线更新，并报告不更新版本。source anchoring、模式/响应目标及是否使用额外采样噪声逐项标注。
\small
\begin{longtable}{L{2.9cm}L{2.4cm}L{6.9cm}}
\toprule
对照 & 目标记录/更新 & 比较目的与定义\\\midrule\endfirsthead
\toprule 对照 & 目标记录/更新 & 比较目的与定义\\\midrule\endhead
Direct source execution & 无校准、无更新 & source prior 经共享目标 adapter 直接执行，测接口落地后的基础差距。\\
Shared prior + context & 相同 $B$；同预算离线更新 & 直接把校准 context/历史提供给共享条件策略，不使用 source-conditioned law transport；检验收益是否仅来自条件信息。\\
Action calibration only & 相同 buffer；仅拟合接口残差 & 校准增益/延迟或受限动作映射，固定源行为结构；目标执行调试也计入 $B$。\\
GAT-style task-space residual & 相同 $B$ 与模拟更新量 & 以当前状态/历史和单步源指令拟合 grounding residual，不加入源--目标响应 coupling 目标；明确是统一接口改编，不冒称原版 GAT\cite{gat}。\\
Reweighting-only & 相同模型与更新预算 & 仅改变已有源候选质量、不移动指令支持；检验支持形变是否必要。\\
Deterministic transport & 相同目标、模型与容量 & source-conditioned 短时 mapping 不输入额外噪声；随机性仍来自固定 source prior。\\
DR + calibration & 相同 source 数据、$B$ 与离线预算 & 在声明机制族中做 DR，并利用同一校准信息/模型更新目标策略；检验普通鲁棒化加辨识能否解释收益\cite{dr}。\\
FORGE-style contact RL & 接触任务；同 $B$ 与模拟预算 & 复用 force-aware DR 与超力成本思想，以统一 13D 接口、相同观测训练；标为 adapted 版本，而非原生 Franka 数字\cite{isaaclabtasks}。\\
Classical contact search & 相同已知模型与 probe buffer & 夹具定位使用 hybrid force/position 控制；插入使用限力搜索/回撤重试。参数只在共同模拟/校准预算内调整，计入所有目标尝试。\\
Full transport & 相同 $B$、模拟量与更新量 & 完整响应校准、任务几何、source-conditioned transport 与可行约束。\\
\bottomrule
\end{longtable}
\normalsize

\textbf{完整系统与诊断分开。} ASID 保留主动探索、辨识及校准后模拟训练流程，作为单列的 end-to-end 系统对照；匹配目标总预算并报告新增模拟量，不能只取其辨识模块代表原方法\cite{asid}。较高计算预算的 motion-planning teacher、已知真值模型规划器及前述可行性 oracle 仅进入诊断图，不与主表合并，也不向任何主方法提供目标示范。预算匹配的 calibrated-model 规划/控制可单列为强执行对照，报告其在线计算与延迟；它与拥有目标真值的 teacher 不同。

LSA、OT-sim2real 与 LAD 如补做公开协议，保留各自所需的随机目标交互、demonstration 或 aligned-action 监督，另列协议 B，不能将额外监督隐藏在校准主表\cite{lsa,otreal,lad}。Human2Sim2Robot 与 TRANSIC 仅作为共同物体语义/人工纠正监督的相关扩展\cite{human2sim,transic}。不新增 latent DAG、causal module discovery 或 selective-update 贡献。

\textbf{资产与工程边界。} Factory/FORGE 的 peg/gear/fastener 和 AutoMate 的 plug/socket 可作为资产及控制参考；官方原生任务主要使用 Franka，跨臂需要新增适配，并重新检查重力、关节限制、全臂/持件碰撞及接触设置\cite{isaaclabtasks}。Task 1--3 场景及 Task 4 的移动身体组合、源多模式、干预和恢复指标由本协议构建，不把现成资产描述为已实现的新 benchmark。MuJoCo Menagerie 的独立 Go2、Z1、G1 等模型可作为候选资产入口，组合安装、工具和 WBC 仍需新增验证\cite{menagerie}。

\textbf{统计计划。} 仿真每个 task--body--controller--mechanism--budget cell 至少 5 个独立训练/校准 seeds，每 seed 至少 50 次冻结后的评估 episodes；所有方法共用预生成 reset/扰动列表，并报告成功率区间和跨 seed/机制层级的不确定性。实物样本量待平台和单次成本确定后预注册，分别报告目标 probe、最终任务评估及 reset 时间。控制频率上的采样点不当作独立试次；不以显著性不足或模式未覆盖为由事后移除任务。

\section{Discussion}
\paragraph{当前结果要求的理论定位。}
P1.5--P1.7 支持将模型校准、执行阶段支持、有限 reset 的选择误差与模式质量
取舍分层；它们没有建立 source anchoring 相对同模型、同容量控制器的任务优势。
尤其 v2 的响应指标下降伴随一个额外实际失败，因此不能把 transport objective
的最优性写成目标任务的优越性。当前可保留的保证是同一声明集合内、条件于
统一误差的任务预算与响应比较，以及不变前缀和支持/停止质量的边界。
若后续具有独立任务意义的结构指标仍不能区分强控制对照，论文应收窄为
可行支持受限的校准控制/行为 law 研究，不能用更大的模型或方法名称补足证据。

论文成立需要同时展示：目标支持/行为 law 的改变；固定接口后 transport 的额外收益；同预算下干预校准的效用；任务成功与模式覆盖、风险的共同改善。

若收益仅来自接口增益校准，则应定位为 action adaptation；若 task-aware ASID 已达到相同效果，则需要谨慎评估 transport 的增量价值；若 ordinary response prediction 与完整方法等效，则缩小 causal 主张。目标分布不必全局复制源分布，必要的模式改变应由可行性和任务结果解释。

\section{Limitations}
\begin{enumerate}
\item 校准与任务共享的结构化机制假设、共同状态充分性及实际访问支持决定泛化范围；未知迟滞/接触分支可能破坏模型。
\item task geometry 与 source anchoring 可能引入偏置；不能把 transport objective 最优等同于任意目标 task-only 最优。
\item 短时模型的组合误差、rare mode 概率与 model-based CVaR 需要独立验证；模型风险约束不是实机安全证书。
\item source-conditioned map/flow 不保证保留所有有效 modes，也不保证创造任意新模式；低支持区域需要新增交互或更宽 prior。
\item intervention 支持规定协议下的响应辨识，不恢复唯一 latent DAG 或任意跨身体个体反事实。
\item 首版仿真采用结构化接触状态，实物平台待确定；不能直接外推到通用视觉、任意软体机器人或真实双臂迁移。
\item Task 4 依赖具备声明末端/力接口的 WBC 和共同源支持；同工位局部操作成功不等于学得导航、步态或任意身体的全身操作策略。身体专属可行性与模型误差需单独验收。
\end{enumerate}

\section{Conclusion}
本文希望通过有限干预校准目标执行机制，将源任务条件行为分布转化为可执行目标行为分布。完整首版应包含 source prior、明确的 action decoding 接口、校准响应、摊销 transport、分布与闭环指标，以及同预算已有方法比较。方法与实验决定贡献，理论解释其依据与边界。

\appendix
\section{Supporting Derivations}
\subsection{Identifiability and Fixed-Design Uncertainty}
若两个参数产生相同校准结果，则 $v=\theta_1-\theta_2\in\ker A_{\mathcal B}$。相关响应一致当且仅当 $B_gv=0$，直接得到主文核空间条件。

固定设计且 $\xi\sim\mathcal N(0,\sigma^2I)$ 时
\begin{equation}
 \hat b_g=B_gA_{\mathcal B}^\dagger y,\qquad
 \Sigma_g=\sigma^2B_g(A_{\mathcal B}^\top A_{\mathcal B})^\dagger B_g^\top.
\end{equation}
在可辨识条件下，以概率至少 $1-\eta$，
\begin{equation}
 \|\hat b_g-b_g\|\le
 \sqrt{\tr\Sigma_g}+\sqrt{2\|\Sigma_g\|_{\rm op}\log(1/\eta)}.
\end{equation}
多个查询需联合控制；真实响应模型还需结构偏差、模式概率与噪声 law 的误差。小均值误差不自动推出整个 kernel 接近。

\subsection{Propagation and Objective Error}
对保持同一源参考的真实和预测增广状态 law 应用 triangle inequality：
\begin{equation}
 d_{k+1}\le
 W_1(\mu_k\mathbb K_t^\varphi,\mu_k\widehat{\mathbb K}_t^\varphi)
 +W_1(\mu_k\widehat{\mathbb K}_t^\varphi,\hat\mu_k\widehat{\mathbb K}_t^\varphi)
 \le\varepsilon_{\mathcal B}+Ld_k.
\end{equation}
对 trajectory cost 必须使用逐步 coupling，而非仅边缘 law；在对应的 kernel coupling 和统一 Lipschitz 条件下，
\begin{equation}
 \varepsilon_F\le
 (L_C+\lambda_sL_c)\varepsilon_{\mathcal B}
 \sum_{k=1}^{H}\sum_{j=0}^{k-1}L^j.
\end{equation}
同一个 source coupling、相同 transport 输入及真实/模型共同可行集合为前提。风险 tightening 不成立时，模型可行解未必真实可行。UOT 增强与不连续事件需单独误差分析。

主文目标界通过加减模型目标获得：
\begin{align}
 F_t(\hat\varphi)-F_t(\varphi^*)
 &= [F_t(\hat\varphi)-\widehat F_t(\hat\varphi)]
 +[\widehat F_t(\hat\varphi)-\widehat F_t(\varphi^*)]
 +[\widehat F_t(\varphi^*)-F_t(\varphi^*)]\\
 &\le2\varepsilon_F+\varepsilon_{\rm opt}.
\end{align}
这里优化的是约定 transport 目标，不移除 source regularization 后重新解释为无约束 RL regret。

\section{Suggested Main-Text Allocation and First Implementation}
以约 10 页研究主文为示意：Introduction 1.1，Related Work 0.6，Preliminary 1.1，Method 3.4，Supporting Analysis 0.7，Experiments 2.5，Discussion/Limitations/Conclusion 0.6。此为写作比例，不是任何 venue 的官方页数规则。完整推导、probe 列表、控制接口、原版/改编 baseline 配置及完整统计进入附录。

执行顺序：先实现并验证新的 13D task-space adapter、共同坐标/单位与全臂/持件碰撞接口；在 Panda 上训练并冻结前三项任务所需的多模态 source prior；先跑持件进入受限工位的 Panda$\to$xArm7 交叉矩阵，再加入夹具定位与 keyed connector 插入恢复。随后完成 $B\in\{0,10,30,100\}$ 的探针、held-out 响应与模式消融；接入 UR5e 后分别评估 frozen-transport context-only 泛化和校准 twin 上的 few-shot adaptation。Task 4 复用冻结的 Task 3 prior，先完成固定/移动 xArm7 的身体耦合对照与轮式 WBC，再接入带臂四足及具备接口的人形配置，按 Section~\ref{sec:task4} 验收同任务扩展。目标真值访问始终限于计数 probes 与封存最终评估。实机型号、控制能力和 F/T 接口确定并验证后再落地相同工位任务；不预先假定 Nova5 可用。首版先比较 deterministic/stochastic residual MLP，conditional flow 与视觉留作模型扩展。每一步检查普通动作校准、共享 context 策略或 calibrated-model 控制能否解释全部收益。

先在 CPU 以小 batch 跑通同一正式数值内核、模型更新与最小闭环，再在正式规模训练前完成 GPU qualification：按机器人/场景静态拓扑分桶，环境/候选/ensemble 批轴向量化，技能与物理时间轴使用 scan；热路径不包含 NumPy 转换或逐样本 host 同步。先在代表性机械臂场景验证 scalar/batched/microbatch 一致性、编译次数、吞吐和显存，再随夹具/插入及移动身体拓扑增量验收。Task 4 的 WBC 若依赖 host QP、外部 controller 或 FFI，必须单列 end-to-end 时延与批量成本，不能只计 transport 前向；未经验证的 host WBC 不作为 GPU collector 热路径。单卡合格后验证多卡数据并行与全局风险归约。详细接口、状态匹配、训练路径、阶段交付与验收见项目文档 \url{docs/development/cross-embodiment-transport-plan.md}。

最重要的图：前三项 3D 工位任务及源模式，Task 4 的同一 connector 在固定/轮式/腿式身体上的操作与耦合响应；源指令 law 与实际响应 law 的差别；方法四模块；同预算任务结果；held-out 响应 discrepancy、固定 $U_s$ 的噪声消融与已验证模式覆盖；机器人$\times$控制器交叉矩阵。主要表：统一接口/预算的受控对照、单列的 ASID 完整系统、跨 DoF 与移动身体的冻结/适应两种协议、任务与分布指标。高计算量 teacher 和目标真值可行性 oracle 只用于封存诊断，不混入主表。source sample count 不作为整篇论文的核心横轴。

\begin{thebibliography}{99}
\small
\bibitem{lsa} T. Wang, D. Bhatt, X. Wang, N. Atanasov. Cross-Embodiment Robot Manipulation Skill Transfer using Latent Space Alignment. arXiv:2406.01968, 2024. \url{https://arxiv.org/abs/2406.01968}
\bibitem{mirage} L. Y. Chen et al. Mirage: Cross-Embodiment Zero-Shot Policy Transfer with Cross-Painting. RSS, 2024. \url{https://robot-mirage.github.io/}
\bibitem{fada} A. Xie et al. FADA: Few-Shot Domain Adaptation via Dynamics Alignment for Humanoid Control. arXiv:2606.28476, 2026. \url{https://arxiv.org/abs/2606.28476}
\bibitem{lad} E. Bauer, E. Nava, R. K. Katzschmann. Latent Action Diffusion for Cross-Embodiment Manipulation. ICRA, 2026; arXiv:2506.14608v4. \url{https://arxiv.org/abs/2506.14608}
\bibitem{xskill} M. Xu et al. XSkill: Cross Embodiment Skill Discovery. CoRL, 2023. \url{https://xskill.cs.columbia.edu/}
\bibitem{human2sim} T. G. W. Lum, O. Y. Lee, K. Liu, J. Bohg. Crossing the Human-Robot Embodiment Gap with Sim-to-Real RL using One Human Demonstration. CoRL, 2025. \url{https://proceedings.mlr.press/v305/lum25a.html}
\bibitem{dr} X. B. Peng et al. Sim-to-Real Transfer of Robotic Control with Dynamics Randomization. ICRA, 2018. \url{https://arxiv.org/abs/1710.06537}
\bibitem{gat} J. P. Hanna, P. Stone. Grounded Action Transformation for Robot Learning in Simulation. AAAI, 2017. \url{https://ojs.aaai.org/index.php/AAAI/article/view/11044}
\bibitem{rma} A. Kumar et al. RMA: Rapid Motor Adaptation for Legged Robots. RSS, 2021. \url{https://arxiv.org/abs/2107.04034}
\bibitem{asid} M. Memmel et al. ASID: Active Exploration for System Identification in Robotic Manipulation. ICLR, 2024. \url{https://weirdlabuw.github.io/asid/}
\bibitem{spi} N. Sobanbabu et al. Sampling-based System Identification with Active Exploration for Legged Sim2Real Learning. CoRL, 2025. \url{https://proceedings.mlr.press/v305/sobanbabu25a.html}
\bibitem{optimalexploration} A. Wagenmaker, G. Shi, K. G. Jamieson. Optimal Exploration for Model-Based RL in Nonlinear Systems. NeurIPS, 2023. \url{https://proceedings.neurips.cc/paper_files/paper/2023/hash/31e018f43ab9c7065c058cc2c5848128-Abstract-Conference.html}
\bibitem{otreal} S. Cheng, L. Ma, Z. Chen, A. Mandlekar, C. Garrett, D. Xu. Generalizable Domain Adaptation for Sim-and-Real Policy Co-Training. NeurIPS, 2025. \url{https://arxiv.org/abs/2509.18631}
\bibitem{densityflow} X. Duan, D. Chen. Trajectory-Optimized Density Control with Flow Matching. arXiv:2510.06666, 2025. \url{https://arxiv.org/abs/2510.06666}
\bibitem{dara} J. Liu, H. Zhang, D. Wang. DARA: Dynamics-Aware Reward Augmentation in Offline Reinforcement Learning. ICLR, 2022. \url{https://arxiv.org/abs/2203.06662}
\bibitem{h2oplus} H. Niu et al. H2O+: An Improved Framework for Hybrid Offline-and-Online RL with Dynamics Gaps. ICRA, 2025. \url{https://arxiv.org/abs/2309.12716}
\bibitem{transic} Y. Jiang et al. TRANSIC: Sim-to-Real Policy Transfer by Learning from Online Correction. CoRL, 2024. \url{https://arxiv.org/abs/2405.10315}
\bibitem{isaaclabtasks} Isaac Lab contributors. Available Environments: Factory, FORGE, and AutoMate. Official documentation, accessed 2026-10-12. \url{https://isaac-sim.github.io/IsaacLab/main/source/overview/environments.html}
\bibitem{umilegs} H. Ha, Y. Gao, Z. Fu, J. Tan, S. Song. UMI-on-Legs: Making Manipulation Policies Mobile with Manipulation-Centric Whole-body Controllers. CoRL, 2024; PMLR 270:5254--5270, 2025. \url{https://proceedings.mlr.press/v270/ha25a.html}
\bibitem{menagerie} MuJoCo Menagerie contributors. MuJoCo Menagerie: Robot model collection. Repository, accessed 2026-10-12. \url{https://github.com/google-deepmind/mujoco_menagerie}
\end{thebibliography}
\end{document}
